\documentclass[11pt]{article}

\usepackage[T1]{fontenc}
\usepackage[utf8]{inputenc}
\usepackage{lmodern}
\usepackage{microtype}
\usepackage[margin=0.95in]{geometry}
\usepackage{amsmath,amssymb,amsthm,mathtools}
\usepackage{array,booktabs,longtable}
\usepackage{xcolor}
\usepackage{graphicx}
\usepackage{placeins}
% Support compilation from the repository root or this source directory.
\graphicspath{{data/runtime/poisson-fixed-participation/independent-batch-sweep-eps-1-2-4-8/figures/}{../../data/runtime/poisson-fixed-participation/independent-batch-sweep-eps-1-2-4-8/figures/}}
\usepackage[colorlinks=true,linkcolor=blue!55!black,urlcolor=blue!55!black]{hyperref}

\setlength{\emergencystretch}{3em}
\setlength{\parskip}{0.45em}
\setlength{\parindent}{0pt}

\newtheorem{theorem}{Theorem}[section]
\newtheorem{proposition}[theorem]{Proposition}
\newtheorem{lemma}[theorem]{Lemma}
\newtheorem{corollary}[theorem]{Corollary}
\theoremstyle{remark}
\newtheorem{remark}[theorem]{Remark}

\newcommand{\R}{\mathbb R}
\newcommand{\N}{\mathbb N}
\newcommand{\E}{\mathbb E}
\newcommand{\Normal}{\mathcal N}
\newcommand{\TV}{\operatorname{TV}}
\newcommand{\AO}{\mathrm{AO}}
\newcolumntype{L}[1]{>{\raggedright\arraybackslash}p{#1}}

\title{Batch Size and Frozen-Gradient Variance\\
at Fixed Expected Participation}
\author{}
\date{}

\begin{document}
\maketitle

\begin{abstract}
We compare batch sizes for equal-weight averages of frozen-gradient estimates
under a fixed expected participation budget and a common privacy target.
For a positive integer $E$, a run with $T\geq E$ independent releases uses
Poisson sampling at rate $q=E/T$. Gaussian noise is calibrated using the
exact two-sided privacy profile of the canonical sampled-Gaussian transcript
under add/remove adjacency.

Let $v_{\mathrm{FB}}$ denote the calibrated privacy variance of the full-batch
average, and set $a=\Delta/\sqrt{v_{\mathrm{FB}}}$, where $\Delta>0$ is the
sensitivity. For $\varepsilon>0$ and $0<\delta<1$, we prove that
\[
  a^2\leq\frac{6E}{3E+2}\,\varepsilon
\]
guarantees strictly smaller variance at full batch than at every integer
horizon $T>E$, for every nonnegative frozen-gradient energy. We also give a
sharper horizon-dependent condition. Both finite guarantees hold when
privacy is calibrated for the canonical average alone. A certified counterexample at
$(\varepsilon,\delta)=(1/4,3/4)$ shows that full-batch optimality need not hold
without a privacy restriction.

For transcript calibration, in the sparse limit $T\to\infty$ with $E$ fixed,
we determine the exact
critical-window privacy profile. For $0<\delta<1-e^{-E}$, we identify the
resulting limiting offset of
$\Delta/\sigma_{E,T}-\sqrt{2\log T}$, where $\sigma_{E,T}$ is the calibrated
raw noise standard deviation. Consequently, full batch eventually has
strictly smaller transcript-calibrated variance at every nonnegative
frozen-gradient energy. Calibration for the canonical average alone permits
Gaussian variance bounded uniformly in $T$.
These results concern frozen-gradient averages, not optimizer trajectories
or final training losses.
\end{abstract}

\section{Introduction}
Increasing the batch size changes both sampling variability and the noise
needed for privacy. We compare these effects while fixing the expected
number of participations of each record. With a positive integer budget $E$,
a horizon $T\geq E$ uses independent Poisson sampling at rate $q=E/T$.
The question is when full batch, $T=E$, minimizes the variance of an unbiased
average under a common privacy target. Optimality is within this family of
independent Poisson releases with homogeneous independent Gaussian noise
and equal-weight averaging of fixed clipped contributions.

Our main finite result gives an explicit sufficient condition under which
full batch strictly minimizes this variance over every finite smaller-batch
horizon (Corollary~\ref{cor:fixed-epoch-all-horizon}). It follows from a
sharper comparison at each fixed horizon
(Theorem~\ref{thm:rate-adaptive-pointwise}). The proof uses a tail event
of the average, reduces it to a binomial count, and establishes a sign inequality for a Gaussian CDF averaged over
that count. It first proves the comparison for average-only calibration;
data processing then gives the transcript comparison. A finite
counterexample then shows why a privacy qualification is necessary
(Theorem~\ref{thm:strict-interior-counterexample}).

The sparse limit gives a complementary conclusion. When $T\to\infty$ with
$Tq=E$ fixed, the calibrated raw noise tends to zero, but the privacy variance
of the average diverges. Theorem~\ref{thm:smallbatch} identifies both the exact
constant offset in the standardized shift and the leading variance growth.
Thus sufficiently small batches are worse than full batch throughout the
stated subcritical privacy regime, even when the finite sufficient condition
does not apply.

When only the canonical average is protected, its calibrated Gaussian
variance stays bounded uniformly in the horizon. This contrasts with the
sparse behavior under transcript calibration.

R\"ais\"a, J\"alk\"o, and Honkela~\cite{raisa2024} study fixed positive
sampling rates as the number of releases grows. Kalinin~\cite{kalinin2024}
proves their sufficient-condition conjecture for the single-release
comparison, including $\varepsilon,q\geq4\delta$. Neither comparison holds
$Tq$ fixed while varying the transcript length. We explain the distinction
in Section~\ref{sec:sparse}. The appendices contain the binomial proof,
the sparse-limit argument, the rational certificate for the counterexample,
and several supplementary comparisons.

\section{Model and calibration}\label{sec:model}
Fix the same parameter value across neighboring datasets and one coordinate
of its clipped per-example gradients. For a dataset of size $N$, let the fixed contributions satisfy
$|\widetilde g_i|\leq\Delta$, where $\Delta>0$, and write
\[
 G=\sum_{i=1}^N\widetilde g_i,\qquad
 S=\sum_{i=1}^N\widetilde g_i^2.
\]
Our target is the clipped full sum $G$. Fix an expected participation budget
$E\in\N_{>0}$ and a privacy target $\varepsilon\geq0$, $0<\delta<1$.
For each integer horizon $T\geq E$, use
\begin{equation}\label{eq:rate}
 q=E/T,\qquad B=qN=EN/T,\qquad Tq=E,\qquad TB=EN.
\end{equation}
Here $B$ is the expected batch size on the dataset being evaluated. The
privacy comparison holds $q$ and the other mechanism parameters fixed across
neighboring datasets; $B=qN$ does not prescribe changing $q$ with dataset
size. Increasing $T$ produces smaller
expected batches while preserving $EN$ expected record participations across
the run. Participation is independent across records and releases;
$E$ is not a deterministic participation count. Full batch is $T=E$, and
its nearest smaller-batch competitor is $T=E+1$.

\paragraph{The estimator and its variance.}
Let $B_{it}\sim\operatorname{Bernoulli}(q)$ be independent inclusion
indicators, and let $\xi_t\sim\Normal(0,\sigma^2)$ be independent of all
indicators and of each other. Before calibrating the raw standard deviation
$\sigma\geq0$, define
\[
 \widehat G_{E,T}=\frac1E\sum_{t=1}^T
       \left(\sum_i B_{it}\widetilde g_i+\xi_t\right).
\]
We hold the dataset and parameter value fixed; the randomness comes from the
inclusion indicators and Gaussian noises. Since $Tq=E$, the estimator satisfies
$\E\widehat G_{E,T}=G$. Independence makes all cross-covariances vanish, so
\begin{equation}\label{eq:uncalibrated-variance}
 \begin{aligned}
 \operatorname{Var}(\widehat G_{E,T})
 &=\frac1{E^2}\sum_{t=1}^T\sum_i
   \widetilde g_i^2\operatorname{Var}(B_{it})
   +\frac1{E^2}\sum_{t=1}^T\operatorname{Var}(\xi_t)\\
 &=\frac{Tq(1-q)}{E^2}S+\frac{T\sigma^2}{E^2}\\
 &=\frac{1-q}{E}S+\frac{T\sigma^2}{E^2}.
 \end{aligned}
\end{equation}
The first term is the Poisson sampling variance, evaluated using
$\operatorname{Var}(B_{it})=q(1-q)$; $S$ is the sum of squared fixed
contributions.
Unbiasedness refers to the clipped full sum, not the unclipped gradient.
The remaining question is how the Gaussian variance required for privacy
changes with the horizon.

Writing $Y_t=\sum_i B_{it}\widetilde g_i+\xi_t$ and
$K_{i,T}=\sum_t B_{it}$ makes the normalization explicit:
\begin{equation}\label{eq:participation-decomposition}
 \widehat G_{E,T}=\frac1T\sum_{t=1}^T\frac{Y_t}{q}
 =\sum_i\frac{K_{i,T}}E\widetilde g_i+\frac1E\sum_{t=1}^T\xi_t.
\end{equation}
Since $\E K_{i,T}=E$, the data term remains unbiased as $T$ increases.
The accumulated raw Gaussian noise is divided by the fixed budget $E$.
When noise is instead specified by a dimensionless multiplier
$\sigma_{\mathrm{DP}}$, the conversion is $\sigma=\Delta\sigma_{\mathrm{DP}}$.

\paragraph{Protected output.}
Transcript calibration protects all individual releases jointly.
Average-only calibration protects their aggregate; it does not guarantee
privacy if the individual releases are also disclosed. The finite dominance
theorem holds for either output choice. The sparse divergence is specific
to transcript calibration, so these are different release choices, not
merely different accountants for the same public output.

\paragraph{The canonical privacy experiment.}
At raw noise standard deviation $\sigma\geq0$, define
\[
 Q_\sigma=\Normal(0,\sigma^2),\qquad
 P_{q,\sigma}=(1-q)\Normal(0,\sigma^2)+q\Normal(\Delta,\sigma^2).
\]
The differing record contributes the fixed shift $\Delta$ when selected;
$\Normal(\mu,0)$ denotes a point mass. We calibrate the noise using the exact
privacy profile of this binary experiment. It is the canonical sampled-Gaussian
comparison used in~\cite{raisa2024}.

For probability laws $P,Q$ and $\alpha\geq0$, write
\[
 H_\alpha(P,Q)=\sup_A\{P(A)-\alpha Q(A)\},
\]
where the supremum is over measurable events. This definition includes
singular zero-noise laws. When $P\ll Q$, it also equals
$\E_Q[(dP/dQ-\alpha)_+]$. Following~\cite{raisa2024}, write the accounting
oracle of the canonical $T$-release experiment as
\begin{equation}\label{eq:canonical-accountant}
 \begin{split}
 \AO_S(\sigma,\Delta,q,T,\varepsilon)=\max\bigl\{&
 H_{e^\varepsilon}(P_{q,\sigma}^{\otimes T},Q_\sigma^{\otimes T}),\\
 &H_{e^\varepsilon}(Q_\sigma^{\otimes T},P_{q,\sigma}^{\otimes T})\bigr\}.
 \end{split}
\end{equation}
The subscript $S$ in $\AO_S$ labels the sampled-Gaussian mechanism; it does
not denote the gradient energy $S$.
For fixed $(\varepsilon,\delta,\Delta)$, $T\in\N_{>0}$, and $q\in[0,1]$,
define the calibrated raw Gaussian noise standard deviation by
\begin{equation}\label{eq:calibration}
 \begin{split}
 \sigma(q,T)&=\inf\{\sigma\geq0:
                  \AO_S(\sigma,\Delta,q,T,\varepsilon)\leq\delta\},\\
 \sigma_{E,T}&=\sigma(E/T,T).
 \end{split}
\end{equation}
We use $\phi$, $\Phi$, and $\overline\Phi=1-\Phi$ for the standard Gaussian
density, CDF, and upper tail, and $\Phi^{-1}$ for its quantile function.

\paragraph{The calibrated variance benchmark.}
With the calibrated noise from \eqref{eq:calibration}, define
\begin{equation}\label{eq:variance}
 V_E(T,S)=\frac{T-E}{ET}S+v_{E,T},
 \qquad v_{E,T}=\frac{\sigma_{E,T}^2}{Tq^2}
                   =\frac{T}{E^2}\sigma_{E,T}^2.
\end{equation}
We call $v_{E,T}$ the \emph{privacy variance} of the average. The noise added
to each raw release has standard deviation $\sigma_{E,T}$, whereas one
unbiased full-sum estimate has noise standard deviation $\sigma_{E,T}/q$.
In the effective-noise notation of~\cite{raisa2024}, for $q>0$,
\begin{equation}\label{eq:effective-noise-normalization}
 \sigma_{\mathrm{eff}}(q,T)=\frac{\sigma(q,T)}q,\qquad
 v_{E,T}=\frac{\sigma_{\mathrm{eff}}(E/T,T)^2}{T}.
\end{equation}
The factor $1/T$ accounts for the additional averaging step. The comparison
uses full-sum normalization rather than dataset-mean normalization. Write
\begin{equation}\label{eq:privacycomponents}
 v_{\mathrm{FB}}=v_{E,E}=\sigma_{E,E}^2/E,
 \qquad v_T=v_{E,T},\quad v_E=v_{\mathrm{FB}}
\end{equation}
when $E$ is fixed.

For the connection to clipped-query privacy, the recordwise query is held
fixed across neighboring datasets, and the clipping bound holds for every
record in the query domain. The canonical experiment then bounds the
corresponding neighboring transcript profiles: a smaller differing
contribution is obtained by a Gaussian channel, and the shared sampled
contributions are common independent randomness. We give the argument in
Appendix~\ref{app:canonical-domination}. A particular neighboring pair may
therefore have a smaller privacy profile than the canonical pair. Noise is
calibrated uniformly using $\Delta$, whereas $S$ describes the particular
frozen dataset on which utility is evaluated. The primary calibration protects
the complete transcript, while utility uses its average. We also consider
direct calibration of the canonical average below.

\paragraph{Calibration in variance units.}
At a candidate privacy variance $v$, put
\begin{equation}\label{eq:runvariance-accountant}
 \mathcal A_{E,T}(v;\varepsilon)
 =\AO_S\!\left(E\sqrt{v/T},\Delta,E/T,T,\varepsilon\right).
\end{equation}
The change of variables $v=T\sigma^2/E^2$ is an increasing continuous
bijection of $[0,\infty)$, so
\[
v_{E,T}=\inf\{v\geq0:\mathcal A_{E,T}(v;\varepsilon)\leq\delta\}.
\]

\paragraph{Calibration for the canonical average alone.}
For comparison, suppose that only the average of the canonical transcript
is released. At Gaussian variance $v\geq0$, its two laws are
\begin{equation}\label{eq:average-experiment}
 \overline P_{E,T,v}
 =\sum_{k=0}^T\Pr\{K=k\}\Normal(k\Delta/E,v),\qquad
 \overline Q_v=\Normal(0,v),\qquad
 K\sim\operatorname{Binomial}(T,E/T).
\end{equation}
Define its profile and calibrated Gaussian variance by
\[
 \overline{\mathcal A}_{E,T}(v;\varepsilon)
 =\max\{H_{e^\varepsilon}(\overline P_{E,T,v},\overline Q_v),
         H_{e^\varepsilon}(\overline Q_v,\overline P_{E,T,v})\},
\]
\[
 \overline v_{E,T}=\inf\{v\geq0:
       \overline{\mathcal A}_{E,T}(v;\varepsilon)\leq\delta\}.
\]
The map $(x_1,\ldots,x_T)\mapsto E^{-1}\sum_t x_t$ sends the canonical
transcript at variance $v$ to this pair. Data processing therefore gives
\begin{equation}\label{eq:average-transcript-order}
 \overline{\mathcal A}_{E,T}(v;\varepsilon)
 \leq\mathcal A_{E,T}(v;\varepsilon),\qquad
 \overline v_{E,T}\leq v_{E,T}.
\end{equation}
Both calibrations use the same sampling-variance term $(1-q)S/E$ when
evaluating the frozen-gradient benchmark. Define the corresponding total
variance under average-only calibration by
\begin{equation}\label{eq:average-total-variance}
 \overline V_E(T,S)=\frac{1-q}{E}S+\overline v_{E,T}.
\end{equation}
Appendix~\ref{app:canonical-domination} gives a direct channel from the
canonical average pair to each neighboring average of the fixed query,
so this calibration also provides a uniform privacy bound for that class.

\begin{lemma}[Calibration properties]\label{lem:calibration}
For every finite admissible horizon and fixed $\varepsilon$, the transcript
accountant value is continuous and nonincreasing in the raw noise, including
at zero. Its feasible set is nonempty, and the calibration infimum is
attained. At a positive minimizing noise, the accountant value equals
$\delta$. The same properties hold for the average-only accountant value
as a function of $v$.

At full batch, $v_{\mathrm{FB}}>0$ is independent of $E$ and is uniquely
specified by
\begin{equation}\label{eq:gaussian-calibration}
 G_\varepsilon(a)=\delta,\qquad a=\Delta/\sqrt{v_{\mathrm{FB}}},\qquad
 G_\varepsilon(a)=\Phi\!\left(\frac a2-\frac\varepsilon a\right)
 -e^\varepsilon\Phi\!\left(-\frac a2-\frac\varepsilon a\right).
\end{equation}
At full batch, $\overline v_{E,E}=v_{\mathrm{FB}}$.
The forward and reverse full-batch profiles coincide, and
$G_\varepsilon'(a)=\phi(a/2-\varepsilon/a)>0$ for $a>0$.
\end{lemma}

The proof is in Appendix~\ref{app:calibration}. The full-batch formula is the
analytical Gaussian calibration of Balle and Wang~\cite{ballewang2018}.
Thus $a$ depends only on $(\varepsilon,\delta)$ and
$v_{\mathrm{FB}}=\Delta^2/a^2$. The sufficient conditions below depend
on the participation schedule and privacy target, but not on the sensitivity
scale $\Delta$.
Monotonicity and attainment also give the implication used throughout:
\begin{equation}\label{eq:calibration-transfer}
 \begin{aligned}
 \mathcal A_{E,T}(v_*;\varepsilon)>\delta
 &\quad\Longrightarrow\quad v_{E,T}>v_*,\\
 \overline{\mathcal A}_{E,T}(v_*;\varepsilon)>\delta
 &\quad\Longrightarrow\quad \overline v_{E,T}>v_*.
 \end{aligned}
\end{equation}
Indeed, $v_{E,T}\leq v_*$ would imply
$\mathcal A_{E,T}(v_*;\varepsilon)
\leq\mathcal A_{E,T}(v_{E,T};\varepsilon)\leq\delta$.
The same argument applies to the average-only profile.

Finally,
\[
 V_E(T,S)-V_E(E,S)=\frac{1-q}{E}S+v_{E,T}-v_{\mathrm{FB}}.
\]
For $T>E$ and $S\geq0$, full batch has no larger total variance exactly when
\begin{equation}\label{eq:threshold}
 S\geq S_{\min}(E,T):=
 \frac{E}{1-q}\max\{v_{\mathrm{FB}}-v_{E,T},0\}.
\end{equation}
Thus strict privacy-variance dominance implies strict total-variance
dominance at every nonnegative energy.

\section{Full-batch optimality}\label{sec:near-full}
The main finite conclusion is that full batch strictly minimizes the variance
over every admissible horizon whenever
\[
 a^2\leq\frac{6E}{3E+2}\,\varepsilon,
 \qquad a=\Delta/\sqrt{v_{\mathrm{FB}}}.
\]
We prove a sharper sufficient condition at each fixed horizon first, then
derive this uniform conclusion in Corollary~\ref{cor:fixed-epoch-all-horizon}.
Define
\begin{equation}\label{eq:rate-adaptive-signal-factor}
 \Lambda_{E,T}=\begin{cases}
  6ET/(3ET+2T-4E),&T>2E,\\
  2,&E<T\leq2E.
 \end{cases}
\end{equation}
Equivalently,
$\Lambda_{E,T}^{-1}=\tfrac12+\max\{c_{E,T},0\}/E$ with
$c_{E,T}=(1-2E/T)/3$; the proof explains the role of $c_{E,T}$.

\begin{theorem}[Horizon-dependent full-batch dominance]
\label{thm:rate-adaptive-pointwise}
Let $E,T\in\N_{>0}$ with $T>E$, $\Delta>0$, $\varepsilon>0$, and
$0<\delta<1$. Put $a=\Delta/\sqrt{v_{\mathrm{FB}}}$, where
$v_{\mathrm{FB}}$ is the calibrated full-batch privacy variance. If
\begin{equation}\label{eq:rate-adaptive-signal-condition}
 a^2\leq\Lambda_{E,T}\varepsilon,
\end{equation}
then
\[
 v_{\mathrm{FB}}<\overline v_{E,T}\leq v_{E,T},
\]
and, for every $S\geq0$,
\[
 V_E(E,S)=\overline V_E(E,S)<\overline V_E(T,S)\leq V_E(T,S).
\]
\end{theorem}

\begin{proof}
By~\eqref{eq:calibration-transfer} and~\eqref{eq:average-transcript-order},
it suffices to show that
$\overline{\mathcal A}_{E,T}(v_{\mathrm{FB}};\varepsilon)>\delta$: some event
for the average must have value larger than $\delta$ at the full-batch
variance. We test the event that is optimal at full batch.

Evaluate the smaller-batch canonical experiment at $v_{\mathrm{FB}}$ and
divide the sum of the raw releases by $E\sqrt{v_{\mathrm{FB}}}$. Without the
differing record (the reference law) this standardized sum is $\Normal(0,1)$.
With it (the source law), conditional on the record's participation count
$K\sim\operatorname{Binomial}(T,E/T)$, it is $\Normal(aK/E,1)$. At full batch
$K=E$, and the likelihood ratio $e^{ay-a^2/2}$ exceeds $e^\varepsilon$
exactly on the ray $y>c_*:=a/2+\varepsilon/a$. This \emph{sum event}
therefore attains the full-batch forward profile:
\begin{equation}\label{eq:fullbatch-event-identity}
 \delta=\overline\Phi(c_*-a)-e^\varepsilon\overline\Phi(c_*).
\end{equation}

Put
\[
 Z=K-E,\qquad h=a/E,\qquad
 \theta=E\left(\frac{\varepsilon}{a^2}-\frac12\right).
\]
Then $c_*-a=h\theta$ and $aK/E-c_*=h(Z-\theta)$. Thus $\theta$ is the signed
offset of the full-batch threshold from the full-batch signal mean, in units
of one participation; Appendix~\ref{app:binomial-signs} calls it the
\emph{center}. The sampled value of the sum event is
$\E[\Phi(h(Z-\theta))]-e^\varepsilon\overline\Phi(c_*)$, and its difference
from the full-batch value~\eqref{eq:fullbatch-event-identity} is
\begin{equation}\label{eq:main-count-gap}
 \Gamma_{E,T}(\theta,h)
 :=\E[\Phi(h(Z-\theta))]-\Phi(-h\theta).
\end{equation}
Consequently,
\begin{equation}\label{eq:event-profile-bound}
\mathcal A_{E,T}(v_{\mathrm{FB}};\varepsilon)
 \geq\overline{\mathcal A}_{E,T}(v_{\mathrm{FB}};\varepsilon)
 \geq\delta+\Gamma_{E,T}(\theta,h).
\end{equation}
A positive gap therefore suffices; the tied symmetric case will be handled
by changing the threshold.

Since $\E Z=0$, Jensen's inequality would give $\Gamma_{E,T}\geq0$ at once
if $z\mapsto\Phi(h(z-\theta))$ were convex on the support of $Z$. It is
convex only for $z<\theta$ and concave beyond, so the sign depends on how the
count distributes mass around the inflection. To see which feature of the
count matters, hold $E,T,\theta$ fixed and let $h\downarrow0$. The linear term
cancels because $\E Z=0$, and expansion over the finite support gives
\begin{equation}\label{eq:gap-small-scale}
 \Gamma_{E,T}(\theta,h)
 =\frac{\phi(0)\operatorname{Var}(Z)}{2}
  \left(\theta-\frac{\E Z^3}{3\E Z^2}\right)h^3+O(h^5).
\end{equation}
In this expansion, increasing the center contributes positively, while a
positive third centered moment contributes negatively. This reflects the
competition between the convex and concave parts of the CDF. For the binomial count,
$\E Z^3/(3\E Z^2)=(1-2E/T)/3=c_{E,T}$, so at a fixed center positivity for
small $h$ requires $\theta\geq c_{E,T}$. The expansion is only a necessary
condition. It says nothing about larger $h$.

Theorem~\ref{thm:rate-adaptive-center} supplies the argument at every scale.
It shows that $\Gamma_{E,T}(\theta,h)>0$ for all $h>0$ whenever
$\theta\geq\max\{c_{E,T},0\}$. The one exception is $T=2E$ with $\theta=0$,
where the gap is zero. In particular, for $T>2E$ the small-scale condition is
also sufficient. Since $\theta=E(\varepsilon/a^2-1/2)$, the reciprocal form
of~\eqref{eq:rate-adaptive-signal-factor} gives
\[
 \theta\geq\max\{c_{E,T},0\}
 \iff\frac{\varepsilon}{a^2}\geq\frac12+\frac{\max\{c_{E,T},0\}}{E}
 =\Lambda_{E,T}^{-1}
 \iff a^2\leq\Lambda_{E,T}\varepsilon.
\]
Thus the average-only profile exceeds $\delta$ except possibly when $T=2E$ and
$a^2=2\varepsilon$.

At that endpoint $c_*=a$, and the symmetry of $K-E$ gives zero gap for the
sum event. This does not settle the profile, which is a supremum over events.
The optimal event for the average is where its likelihood ratio $L$ exceeds
$e^\varepsilon$, and the ray $(a,\infty)$ is not of that form near its
endpoint: completing the square gives
\[
 L(a)=e^\varepsilon\E\exp\!\left[
       -\frac{a^2}{2}(K/E-1)^2\right]<e^\varepsilon.
\]
Points just above $a$ therefore lower the event value, and removing them
raises it strictly above $\delta$. Lemma~\ref{lem:strict-average-endpoint}
makes this precise.

In every case the average-only profile exceeds $\delta$ at $v_{\mathrm{FB}}$.
Equations~\eqref{eq:calibration-transfer} and~\eqref{eq:average-transcript-order}
give $v_{\mathrm{FB}}<\overline v_{E,T}\leq v_{E,T}$. Adding the nonnegative
sampling term proves the total-variance comparison for transcript calibration
and also for calibration of the canonical average alone.
\end{proof}

\begin{corollary}[Full-batch dominance at every finite horizon]
\label{cor:fixed-epoch-all-horizon}
Fix $E\in\N_{>0}$, $\Delta>0$, $\varepsilon>0$, and $0<\delta<1$, and put
$a=\Delta/\sqrt{v_{\mathrm{FB}}}$. If
\begin{equation}\label{eq:epoch-uniform-condition}
 a^2\leq\Lambda_E\varepsilon,
 \qquad \Lambda_E:=\frac{6E}{3E+2},
\end{equation}
then, for every finite integer $T>E$,
\[
 v_{\mathrm{FB}}<\overline v_{E,T}\leq v_{E,T},
\]
and, for every $S\geq0$,
\[
 V_E(E,S)=\overline V_E(E,S)<\overline V_E(T,S)\leq V_E(T,S).
\]
\end{corollary}

\begin{proof}
For $T>2E$,
\[
 \Lambda_{E,T}=\frac{6E}{3E+2-4E/T}>\Lambda_E.
\]
For $E<T\leq2E$, $\Lambda_{E,T}=2>\Lambda_E$. The conclusion follows from
Theorem~\ref{thm:rate-adaptive-pointwise}.
\end{proof}

The uniform condition is equivalent to $\theta\geq1/3$ in the proof above.
Strict monotonicity of $G_\varepsilon$ also expresses it directly in terms
of the privacy target:
\begin{equation}\label{eq:epoch-uniform-delta-threshold}
 \delta\leq G_\varepsilon(\sqrt{\Lambda_E\varepsilon})
 \quad\Longleftrightarrow\quad a^2\leq\Lambda_E\varepsilon.
\end{equation}
At $(\varepsilon,\delta)=(1,10^{-5})$, Gaussian calibration gives
$a^2\approx0.0718514$. The condition holds for every positive integer $E$,
since $\Lambda_E\geq6/5$. This measures slack in the sufficient condition,
not a variance ratio.

Appendix~\ref{app:elementary} gives shorter proofs based on full-support
convexity and finite-support secants. Their conditions are more restrictive
than~\eqref{eq:rate-adaptive-signal-condition}; they provide alternative
arguments rather than additional parameter coverage.

\paragraph{Numerical illustration.}
Figure~\ref{fig:fixed-participation-variance} illustrates transcript calibration
at $\varepsilon\in\{1,2,4,8\}$. The normalization is
\begin{equation}\label{eq:plotted-variance-ratio}
 R_E(q):=\frac{v_{E,T}}{v_{\mathrm{FB}}}
 =\frac1q\left(\frac{\sigma_{E,T}}{\sigma_{E,E}}\right)^2,
 \qquad q=E/T,\quad T\in\{E,E+1,\ldots\}.
\end{equation}
Thus reducing raw noise improves the average's privacy variance only if
its ratio to full-batch raw noise falls below $\sqrt q$. A ratio $R_E(q)>1$
implies larger total variance than full batch for every $S\geq0$.
For these four targets at $\delta=10^{-5}$, analytical Gaussian calibration
gives $a^2/\varepsilon<0.347<6/5\leq\Lambda_E$ for every positive integer
$E$. Thus the uniform dominance condition applies throughout the figure:
the sweep illustrates the magnitude of the increase, rather than establishing
its sign. Appendix~\ref{app:numerical-methods} describes the two-sided
profile computation, accuracy checks, and the separate numerical
supplement~\cite{numerical-supplement} containing the code, data, and figure.
The extreme endpoint is a mathematical stress test: $E=1000$ and $q=10^{-6}$
require $T=10^9$ releases.

\begin{figure}[htbp]
\centering
\includegraphics[width=\linewidth]{paper_raw_noise_row.pdf}\\[0.5em]
\includegraphics[width=\linewidth]{paper_variance_row.pdf}
\caption{Transcript-calibrated raw-noise ratio (top) and privacy-variance
ratio $R_E(q)$ (bottom) at fixed expected participation $Tq=E$, with
$\delta=10^{-5}$ and $\Delta=1$. Both axes are logarithmic; $q$ increases
from $10^{-6}$ to full batch at $q=1$. The upper dashed curve $\sqrt q$
is the raw-noise reduction required for equal privacy variance; the lower
dashed line is $R_E=1$. Each $\varepsilon$ panel uses its own calibrated
full-batch denominator, so cross-panel heights compare relative ratios,
not absolute noise or variance. Markers are numerical transcript calibrations
at integer horizons, not certified bounds; lines are guides, not a
monotonicity claim. These are not average-only calibrations.}
\label{fig:fixed-participation-variance}
\end{figure}
\FloatBarrier

\section{Why a privacy condition is necessary}
At zero noise, the privacy profile is $1-(1-E/T)^T$, which tends to
$1-e^{-E}$ as $T\to\infty$; see Appendix~\ref{app:zero-noise}.
The counterexample below has $\delta=3/4<1-e^{-2}$, below this limiting
nonparticipation boundary. Its zero-noise profiles are $1$ at $T=2$ and
$26/27$ at $T=3$, so both calibrated noise levels are positive by continuity.
It therefore shows that a smaller batch can have lower variance even when
neither experiment admits zero noise.

\begin{theorem}[A finite reversal below the nonparticipation boundary]
\label{thm:strict-interior-counterexample}
Let $\Delta=1$, $E=2$, $\varepsilon=1/4$, and $\delta=3/4$. At the adjacent
smaller-batch horizon $T=3$, with $q=2/3$, the calibrated privacy variances
satisfy
\[
 v_3\leq\frac5{32}<\frac4{25}<v_2.
\]
Consequently,
\[
 0\leq S\leq\frac9{400}
 \quad\Longrightarrow\quad V_2(3,S)<V_2(2,S).
\]
\end{theorem}

\begin{proof}
The rational bounds in Appendix~\ref{app:rational-calculation} show that
both directed sampled profiles at $v=5/32$ are below $3/4$, whereas the
full-batch profile at $v=4/25$ exceeds $19/25>3/4$. Feasibility and
\eqref{eq:calibration-transfer} give
$v_3\leq5/32<4/25<v_2$. The sampling term at $E=2,T=3$ is $S/6$, so
\[
 V_2(3,S)\leq\frac S6+\frac5{32}\leq\frac4{25}<v_2=V_2(2,S)
 \qquad(0\leq S\leq9/400).
\]
The endpoint follows from $6(4/25-5/32)=9/400$.
\end{proof}

The same energy interval gives a reversal under average-only calibration:
$\overline V_2(3,S)\leq V_2(3,S)<V_2(2,S)=\overline V_2(2,S)$ for
$0\leq S\leq9/400$.

The example refutes unconditional full-batch optimality: moving from full
batch $T=E$ to its nearest smaller-batch competitor $T=E+1$ can reduce the
variance, even below the nonparticipation boundary.
The stated energy interval is not maximal: under transcript calibration,
the exact reversal condition is
$S<6(v_2-v_3)$, and $9/400<6(v_2-v_3)$. The loose target $\delta=3/4$ shows
that some privacy restriction is necessary; it does not indicate how often
reversals occur at stringent privacy targets.

\section{The sparse regime}\label{sec:sparse}
For fixed $E$, the path $q=E/T\to0$ behaves differently from the fixed-rate
limit. Assume
\begin{equation}\label{eq:interior}
 \Delta>0,\qquad\varepsilon\geq0,\qquad0<\delta<1-e^{-E}.
\end{equation}
\paragraph{Participation does not vanish.}
At zero noise, the differing record is invisible exactly when it is omitted
from every release, so the privacy profile is
$1-(1-E/T)^T\to1-e^{-E}$; see Appendix~\ref{app:zero-noise}. Decreasing $q$
makes participation rare in each release, but the number of opportunities
grows so that $Tq=E$. The probability of at least one participation therefore
tends to $1-e^{-E}$, not to zero, and a target
$\delta<1-e^{-E}$ cannot be met by ignoring participations. In
\eqref{eq:participation-decomposition}, the data term stays unbiased and its
sampling variance tends to $S/E$. The Gaussian term has variance
$T\sigma_{E,T}^2/E^2$, so a bounded privacy variance would require
per-release noise variance $O(1/T)$. We show that the calibrated noise
variance is instead of order $1/\log T$.

\paragraph{The maximum test.}
A participation shifts one of $T$ coordinates by $\Delta$. The noise must hide
this shift among $T$ independent reference coordinates. The simplest test,
which looks at the largest coordinate, already rules out bounded privacy
variance. Suppose that $v_{E,T}\leq M$ along an unbounded sequence of
horizons, with $M>0$. Write $\sigma_T=\sigma_{E,T}$ along that sequence.
These calibrated noises are feasible by Lemma~\ref{lem:calibration} and
satisfy $\sigma_T^2\leq ME^2/T$. For the canonical raw transcript $X$,
consider
\[
 M_T=\{\max_{1\leq t\leq T}X_t>\Delta/2\},\qquad
 u_T=\overline\Phi(\Delta/(2\sigma_T)),
\]
with $u_T=0$ if $\sigma_T=0$. A Gaussian tail bound gives
\[
 Q_{\sigma_T}^{\otimes T}(M_T)\leq Tu_T
 \leq T\exp\!\left(-\frac{\Delta^2T}{8ME^2}\right)\longrightarrow0.
\]
Under the source law, a coordinate exceeds $\Delta/2$ with probability
$p_T=(1-q)u_T+q(1-u_T)$. Hence $Tp_T\to E$, and independence yields
\[
 P_{q,\sigma_T}^{\otimes T}(M_T)=1-(1-p_T)^T\longrightarrow1-e^{-E}.
\]
The forward event value tends to $1-e^{-E}>\delta$, contradicting
feasibility. Every bounded subsequence is therefore excluded, and
$v_{E,T}\to\infty$. This proves divergence, but not its rate.

\paragraph{The critical window.}
To locate the critical window, let $\sigma_T$ now be any noise sequence with
$\Delta/\sigma_T=b_T+s+o(1)$, where $b_T=\sqrt{2\log T}$, and use a maximum
test with threshold $\sigma_T b_T$ rather than the fixed threshold $\Delta/2$
above. Under the reference law its probability is at most
$T\overline\Phi(b_T)\to0$. A selected coordinate exceeds this threshold with
probability tending to $\Phi(s)$.
Independence and $Tq=E$ make the number of selected exceedances converge to a
Poisson variable with mean $E\Phi(s)$; the probability of any exceedance from
unselected coordinates tends to zero. The forward value of this event thus
tends to $1-e^{-E\Phi(s)}$, giving an asymptotic lower bound for the two-sided
profile. Other transcript events might give a larger value. The
likelihood-ratio argument in Appendix~\ref{app:sparse-limit} proves the
matching upper bound and hence the exact limit. For fixed $E$, $\Delta>0$,
$\varepsilon\geq0$, and eventually positive $\sigma_T$,
\begin{equation}\label{eq:main-profile-limit}
 \frac{\Delta}{\sigma_T}-\sqrt{2\log T}\longrightarrow s\in\R
 \quad\Longrightarrow\quad
 \AO_S(\sigma_T,\Delta,E/T,T,\varepsilon)
 \longrightarrow1-e^{-E\Phi(s)}.
\end{equation}
This is the fixed-$E$ case of Theorem~\ref{thm:accountant-limit}.

The threshold $e^\varepsilon$ does not appear in the limit. Let $R_T$ be the
likelihood ratio of the source to the reference transcript law. Under the
reference law, $R_T\to c_s=e^{-E\Phi(s)}<1$ in probability, although
$\E R_T=1$: rare large values, consistent with the visible-selection picture,
carry the missing mean. The directed profiles
$1-\E\min(R_T,e^\varepsilon)$ and $1-e^\varepsilon\E\min(R_T,e^{-\varepsilon})$
involve only bounded functions of $R_T$. Since $c_s<1\leq e^\varepsilon$, they
tend to $1-c_s$ and $(1-e^\varepsilon c_s)_+\leq1-c_s$ for every fixed
$\varepsilon\geq0$.

Inverting~\eqref{eq:main-profile-limit} at $\delta$ gives the calibrated
offset. Let $s_{E,\delta}$ solve $1-e^{-E\Phi(s)}=\delta$. The limit increases
with $s$. So for each $\eta>0$, the candidate noise
$\Delta/(\sqrt{2\log T}+s_{E,\delta}+\eta)$ is eventually infeasible and
$\Delta/(\sqrt{2\log T}+s_{E,\delta}-\eta)$ is eventually feasible.
Monotonicity of the finite accountant traps the calibrated noise between the
two candidates.

\begin{theorem}[Sparse calibration and variance]\label{thm:smallbatch}
Fix a positive integer $E$ and assume~\eqref{eq:interior}. As $T\to\infty$
through integers,
\[
 \frac{\Delta}{\sigma_{E,T}}-\sqrt{2\log T}
 \longrightarrow s_{E,\delta}:=
 \Phi^{-1}\!\left(-\frac{\log(1-\delta)}{E}\right).
\]
In particular,
\begin{equation}\label{eq:sigmaasymptotic}
 \sigma_{E,T}\sim\frac{\Delta}{\sqrt{2\log T}}\longrightarrow0,
\end{equation}
whereas
\begin{equation}\label{eq:privacydivergence}
 v_{E,T}\sim\frac{\Delta^2T}{2E^2\log T}\longrightarrow+\infty.
\end{equation}
For every fixed $S\geq0$,
\begin{equation}\label{eq:total-equivalent}
 V_E(T,S)\sim v_{E,T}.
\end{equation}
\end{theorem}

Appendix~\ref{app:sparse-limit} proves the more general case $Tq_T\to\rho>0$
and then derives this theorem.

\begin{corollary}[Uniform eventual full-batch dominance]\label{cor:eventual}
Under~\eqref{eq:interior}, there is an integer $T_0\geq E$ such that
$V_E(E,S)<V_E(T,S)$ for every $T\geq T_0$ and every $S\geq0$.
\end{corollary}

\begin{proof}
Choose $T_0$ so that $v_{E,T}>v_{\mathrm{FB}}$ for every $T\geq T_0$.
The sampling term is nonnegative, so the same cutoff works for all $S\geq0$.
\end{proof}

This compares each sufficiently sparse horizon with full batch; it does not
assert monotonicity between smaller-batch horizons.

\paragraph{Contrast with average-only calibration.}
\begin{proposition}[Uniform bound for average-only calibration]
\label{prop:average-uniform-bound}
For every integer $T\geq E>0$, $\Delta>0$, $\varepsilon\geq0$, and
$0<\delta<1$,
\[
 \overline v_{E,T}\leq\frac{\Delta^2}{2\pi\delta^2}.
\]
\end{proposition}

\begin{proof}
Use $\TV(P,Q)=\sup_A|P(A)-Q(A)|$. For $m\geq0$ and $v>0$, the Gaussian
density crossing at $m/2$ gives
\[
 \TV(\Normal(m,v),\Normal(0,v))
 =2\Phi\!\left(\frac{m}{2\sqrt v}\right)-1
 \leq\frac{m}{\sqrt{2\pi v}}.
\]
The inequality follows by bounding the standard Gaussian density by
$\phi(0)$. Convexity under mixtures and $\E K=E$ now give
\[
 \TV(\overline P_{E,T,v},\overline Q_v)
 \leq\E\!\left[\frac{K\Delta/E}{\sqrt{2\pi v}}\right]
 =\frac{\Delta}{\sqrt{2\pi v}}.
\]
For $\varepsilon\geq0$, both directed hockey-stick profiles are bounded
by total variation. Thus $v=\Delta^2/(2\pi\delta^2)$ is feasible.
\end{proof}

This bound is a uniform-in-horizon feasibility estimate; it need not be a
close approximation to the calibrated variance.

For fixed $S\geq0$, the sampling term is at most $S/E$, so
\[
 \overline V_E(T,S)\leq\frac SE+\frac{\Delta^2}{2\pi\delta^2}
 \qquad(T\geq E).
\]
Consequently, in the subcritical regime \eqref{eq:interior},
\[
 V_E(T,S)\longrightarrow\infty,\qquad
 \sup_{T\geq E}\overline V_E(T,S)
 \leq\frac SE+\frac{\Delta^2}{2\pi\delta^2}<\infty.
\]
The two release choices differ in what they reveal. At positive noise the
product likelihood ratio of the canonical transcript depends only on the
multiset of observations, so the order of the releases carries no
information. Their individual magnitudes do, and the maximum test exploits
exactly this; averaging discards it. The finite sufficient conditions hold
for both outputs, but the divergence is specific to transcript calibration.

\paragraph{Comparison with fixed-rate asymptotics.}
With $\sigma(q,T)$ defined in~\eqref{eq:calibration}, for fixed $q>0$,
Corollary~5.4 of R\"ais\"a, J\"alk\"o, and
Honkela~\cite{raisa2024} gives
\[
 \frac{\sigma(q,T)}{q\sigma(1,T)}\to1,\qquad
 \sigma(1,T)=\kappa\sqrt T,\qquad
 \kappa=\sqrt{v_{\mathrm{FB}}}.
\]
Hence
\begin{equation}\label{eq:fixedratecontrast}
 \frac{\sigma(q,T)^2}{q^2T}\to v_{\mathrm{FB}}
 \qquad(q>0\text{ fixed}),
\end{equation}
whereas
\begin{equation}\label{eq:fixedepochcontrast}
 \frac{\sigma_{E,T}^2}{q^2T}
 \sim\frac{\Delta^2T}{2E^2\log T}\to\infty
 \qquad(Tq=E).
\end{equation}
These paths impose different participation budgets. The fixed-rate limit
cannot be applied to the vanishing-rate path without a further uniformity
argument.

Ertan and van Dijk~\cite{ertanvandijk2026} obtain the same critical noise scale
through a one-epoch shuffling analysis and a Poisson transfer under zero-out
adjacency. Appendix~\ref{app:separation} relates their separation metric to
the canonical experiment and gives its exact critical-window limit.

\section{Discussion}
The finite theorem identifies a sufficient domain for full-batch optimality;
the reversal shows why a qualification is needed. Neither result classifies
all privacy parameters or proves monotonicity between every pair of
smaller-batch horizons.

For a dataset of size $N$, the expected batch size is $EN/T$. At fixed $N$
it eventually falls below one along the sparse path. Maintaining a
nonvanishing batch size requires a joint limit with $N\to\infty$. The sparse
theorem also fixes $\delta$ and does not address joint scalings such as
$\delta\asymp1/N$. The finite theorem applies pointwise to $\delta=1/N$
whenever its displayed target condition holds.

The variance benchmark assumes independent Poisson participation,
homogeneous independent Gaussian noise, and an equal-weight average of
frozen gradients. The average-only comparison concerns the explicit
canonical pair in~\eqref{eq:average-experiment}. These results do not cover
shuffled or balanced participation, heterogeneous or correlated noise,
adaptive optimizer trajectories, or excess risk.

\paragraph{Formalization and reproducibility.}
The versioned Lean supplement~\cite{rational-supplement} supports the
transcript conclusions of Corollary~\ref{cor:fixed-epoch-all-horizon} and
Theorem~\ref{thm:strict-interior-counterexample}, and those of
Theorem~\ref{thm:rate-adaptive-pointwise} except at $T=2E$,
$a^2=2\varepsilon$. It also supports the general sparse profile and
calibration limits, Theorems~\ref{thm:accountant-limit}
and~\ref{thm:fixed-participation-calibration}, and the fixed-epoch variance
limits and eventual dominance. A current coverage addendum~\cite{formal-coverage}
records exact statements, hypotheses and the translations to this article.
No formal-verification claim is made for the average-only extensions, the
closed symmetric endpoint, continuity at zero and attainment in
Lemma~\ref{lem:calibration}, or the fixed-query domination arguments in
Appendix~\ref{app:canonical-domination}. These are proved conventionally here.

\clearpage
\appendix
\section{Calibration and the canonical experiment}\label{app:calibration}
\begin{proof}[Proof of Lemma~\ref{lem:calibration}]
We establish the finite-noise properties for arbitrary $T\geq1$ and
$0<q\leq1$, then specialize to full batch.

\paragraph{Monotonicity and feasibility.}
If $0\leq\sigma_1\leq\sigma_2$, adding independent centered Gaussian noise of
variance $\sigma_2^2-\sigma_1^2$ to each coordinate maps both transcript laws
at $\sigma_1$ to their counterparts at $\sigma_2$. Data processing for
hockey-stick divergence therefore makes the two-sided profile nonincreasing
in $\sigma$.

The sampled pair is also a common randomized postprocessing of the ordinary
Gaussian pair $\Normal(\Delta,\sigma^2)$ versus $\Normal(0,\sigma^2)$: retain
the input with probability $q$, and otherwise replace it with a fresh
$\Normal(0,\sigma^2)$ draw. Applying this channel independently gives
\[
 \AO_S(\sigma,\Delta,q,T,\varepsilon)
 \leq
 \AO_S(\sigma,\Delta,1,T,\varepsilon).
\]
The Gaussian product profile tends to zero as $\sigma\to\infty$, so the
feasible set is nonempty for every $\delta>0$.

\paragraph{Continuity and attainment.}
For $\sigma>0$, standardize the reference law to $T$ independent standard
normals and write
\[
 R_\sigma(z)=\prod_{i=1}^T
 \left[1-q+q\exp\left(\frac{\Delta z_i}{\sigma}
                         -\frac{\Delta^2}{2\sigma^2}\right)\right].
\]
With $\alpha=e^\varepsilon$, the directed profiles are
\begin{align*}
 H_\alpha(P_{q,\sigma}^{\otimes T}, Q_\sigma^{\otimes T})
 &=1-\E\min(R_\sigma,\alpha),\\
 H_\alpha(Q_\sigma^{\otimes T}, P_{q,\sigma}^{\otimes T})
 &=1-\alpha\E\min(R_\sigma,\alpha^{-1}).
\end{align*}
For fixed $z$, $R_\sigma(z)$ is continuous at positive $\sigma$. The truncated
integrands are bounded independently of $\sigma$, so bounded convergence
proves continuity there.

As $\sigma\downarrow0$, $R_\sigma(z)\to w:=(1-q)^T$ for every fixed $z$.
The same argument gives limits $1-w$ and $(1-\alpha w)_+$. At zero noise,
$Q_0^{\otimes T}$ is concentrated at the origin, and $P_{q,0}^{\otimes T}$
assigns it mass $w$. Direct optimization over events gives precisely these
two profiles. Their maximum is therefore continuous at zero as well.

The feasible set is closed and nonempty. A sequence of feasible noises
converging to its nonnegative infimum shows that the infimum is feasible.
If a positive minimum had profile strictly below $\delta$, continuity would
make a slightly smaller noise feasible. Hence the profile equals $\delta$
at every positive minimizing noise.

\paragraph{Full-batch calibration.}
At privacy variance $v>0$, each of the $E$ full-batch releases has variance
$Ev$. The product likelihood ratio is
\[
 \frac{dP^{\otimes E}}{dQ^{\otimes E}}(x)
 =\exp\left(\frac{\Delta}{Ev}\sum_{t=1}^E x_t
                     -\frac{\Delta^2}{2v}\right).
\]
It depends only on the sum, so the sum preserves both directed profiles.
After division by $E$, the two laws are $\Normal(\Delta,v)$ and
$\Normal(0,v)$. Their directed profiles agree by reflection, and the Gaussian
likelihood-threshold formula gives~\eqref{eq:gaussian-calibration}.

The identity
\[
 e^\varepsilon\phi(-a/2-\varepsilon/a)=\phi(a/2-\varepsilon/a)
\]
gives $G_\varepsilon'(a)=\phi(a/2-\varepsilon/a)>0$. Moreover,
$G_\varepsilon(a)\to0$ as $a\downarrow0$ and
$G_\varepsilon(a)\to1$ as $a\to\infty$, including when $\varepsilon=0$.
Thus the profile is continuous and strictly decreasing in $v>0$, from one
to zero. There is a unique positive $v_{\mathrm{FB}}$ with profile $\delta$,
and it is independent of $E$. Equivalently,
\begin{equation}\label{eq:fullbatch}
 \sigma_{E,E}=\kappa_{\varepsilon,\delta,\Delta}\sqrt E,
 \qquad V_E(E,S)=\kappa_{\varepsilon,\delta,\Delta}^2,
\end{equation}
where $\kappa_{\varepsilon,\delta,\Delta}$ is the one-release Gaussian
calibrated standard deviation. This is the analytical Gaussian calibration
of Balle and Wang~\cite{ballewang2018}.

\paragraph{Average-only calibration.}
For fixed $E,T$, write $p_k=\Pr\{K=k\}$ and $m_k=k\Delta/E$ in
\eqref{eq:average-experiment}. Adding independent Gaussian noise makes the
average-only profile nonincreasing in $v$, and
\eqref{eq:average-transcript-order} gives finite feasibility. At $v>0$,
standardizing the reference observation to $Z\sim\Normal(0,1)$ gives the
likelihood ratio
\[
 \overline R_v(z)=\sum_{k=0}^T p_k
       \exp\!\left(\frac{m_kz}{\sqrt v}-\frac{m_k^2}{2v}\right).
\]
This finite sum is pointwise continuous for $v>0$. As $v\downarrow0$, it
converges pointwise to $p_0=(1-E/T)^T$: every term with $k>0$ tends to zero.
The same bounded-truncation formulas used above prove continuity of both
directed profiles, including at zero. At zero the source has mass $p_0$
at the reference atom and the remaining mass at positive locations, giving
the same directed endpoint values $1-p_0$ and $(1-e^\varepsilon p_0)_+$.
Closedness and the same minimum argument prove attainment and equality to
the target at a positive minimizer. At $T=E$, $K=E$ almost surely, so the
average pair is $\Normal(\Delta,v)$ versus $\Normal(0,v)$ and
$\overline v_{E,E}=v_{\mathrm{FB}}$.
\end{proof}

\subsection{Domination for fixed recordwise queries}\label{app:canonical-domination}
For completeness, consider a fixed recordwise query whose contributions
have magnitude at most $\Delta$. Under add/remove adjacency, let $d$ be the
differing contribution, so $|d|\leq\Delta$. The Gaussian channel
\[
 x\longmapsto\frac d\Delta x+\zeta,
 \qquad
 \zeta\sim\Normal\left(0,\sigma^2\left(1-\frac{d^2}{\Delta^2}\right)\right)
\]
maps $\Normal(0,\sigma^2)$ to itself and $\Normal(\Delta,\sigma^2)$ to
$\Normal(d,\sigma^2)$. It therefore maps the canonical one-release pair to
the pair with differing contribution $d$. The statement also holds at
$\sigma=0$, with degenerate Gaussian noise.

Apply this channel independently across releases, then add the vector of
shared sampled contributions. That vector is independent of the differing
record's inclusion indicators and of the Gaussian noise, and it has the
same law on the two neighboring datasets. These common channels produce
the clipped-sum transcript, so data processing bounds its profile by the
canonical profile. The query must be held fixed across the neighboring
comparison; a dataset-dependent choice of the query would require a
separate argument.

For the average-only comparison, let
$C_T=E^{-1}\sum_{t=1}^T\sum_{i\,\mathrm{shared}} B_{it}\widetilde g_i$
be the shared sampled contribution to the average. It is independent of
the differing record's participation count and of the added noise. Apply
directly to the canonical average the common channel
\[
 X\longmapsto\frac d\Delta X+\zeta+C_T,\qquad
 \zeta\sim\Normal\!\left(0,v\left(1-\frac{d^2}{\Delta^2}\right)\right),
\]
where $\zeta$ and $C_T$ are independent of $X$ and of each other. Under the
reference law the output is $C_T+\Normal(0,v)$. Conditional on $K=k$ under
the source law, it is $C_T+kd/E+\Normal(0,v)$: scaling the input noise and
adding $\zeta$ leaves Gaussian variance $v$. These are exactly the two
neighboring averages when each raw release receives Gaussian noise of
variance $E^2v/T$. The construction
also holds at $v=0$. Data processing therefore bounds each such neighboring
average profile by $\overline{\mathcal A}_{E,T}(v;\varepsilon)$.

\section{Binomial sign comparisons}\label{app:binomial-signs}
Throughout this appendix, $T>E>0$ are integers,
$p=E/T$, $r=1-p$, $K\sim\operatorname{Binomial}(T,p)$, and $Z=K-E$.
The gap $\Gamma_{E,T}$ is defined in~\eqref{eq:main-count-gap}. Its relevant
moments are
\begin{equation}\label{eq:binomial-moments}
 \E Z=0,\qquad \E Z^2=Er,\qquad \E Z^3=Er(1-2p).
\end{equation}
These follow by expanding the sum of $T$ independent centered Bernoulli
variables: every mixed term containing a singleton centered factor vanishes.
Both sign arguments below also use the size-bias identities: if
$K\sim\operatorname{Binomial}(T,p)$ and $L\sim\operatorname{Binomial}(T-1,p)$,
then for every integer $k$,
\begin{equation}\label{eq:binomial-size-bias}
 k\Pr\{K=k\}=Tp\Pr\{L=k-1\},\qquad
 (T-k)\Pr\{K=k\}=Tr\Pr\{L=k\}.
\end{equation}
Adding them gives
$\Pr\{K=k\}=r\Pr\{L=k\}+p\Pr\{L=k-1\}$.

\begin{theorem}[The binomial center comparison]\label{thm:rate-adaptive-center}
Let $T>E>0$ be integers and $h>0$. Then
$\Gamma_{E,T}(\theta,h)>0$ under the following conditions:
\begin{equation}\label{eq:rate-adaptive-center-condition}
 \begin{array}{c|c}
 T>2E & \theta\geq c_{E,T}:=(1-2E/T)/3,\\
 T=2E & \theta>0,\\
 E<T<2E & \theta\geq0.
 \end{array}
\end{equation}
At $T=2E$ and $\theta=0$, the gap is zero.
\end{theorem}

We first show that a nonnegative gap becomes positive when the center is
increased. We then establish the starting comparison at $c_{E,T}$ in the
low-rate case and at zero in the high-rate case. Symmetry handles the
half-rate case.

\subsection*{Propagation in the center}
\begin{lemma}[Propagation of positivity]\label{lem:binomial-propagation}
Fix $h>0$ and $0\leq\theta_1<\theta_2$. If
$\Gamma_{E,T}(\theta_1,h)\geq0$, then
$\Gamma_{E,T}(\theta_2,h)>0$.
\end{lemma}

\begin{proof}
For $t\in\R$, put $k_\theta(t)=h\phi(h(t-\theta))$. For every integer $z$,
\[
 \Phi(h(z-\theta))-\Phi(-h\theta)=\int_0^z k_\theta(t)\,dt.
\]
Let $A_\theta$ be the expectation of the positive integrals and $B_\theta$
the expectation of the absolute values of the negative integrals. Then
$\Gamma_{E,T}(\theta,h)=A_\theta-B_\theta$, and
\[
 \frac{k_{\theta_2}(t)}{k_{\theta_1}(t)}
 =\chi\exp\bigl(h^2(\theta_2-\theta_1)t\bigr),
 \qquad
 \chi=\exp\left(-\frac{h^2}{2}(\theta_2^2-\theta_1^2)\right)>0.
\]
The ratio is strictly larger than $\chi$ for $t>0$ and strictly smaller for
$t<0$. The centered nondegenerate binomial law has positive mass on both
sides of zero, so
\[
 A_{\theta_2}>\chi A_{\theta_1},\qquad
 B_{\theta_2}<\chi B_{\theta_1}.
\]
Subtracting yields
$\Gamma_{E,T}(\theta_2,h)>\chi\Gamma_{E,T}(\theta_1,h)\geq0$.
\end{proof}

\subsection*{Low rate: the critical center}
Assume $T>2E$ and abbreviate $c=c_{E,T}=(1-2p)/3$, so $0<c<1/3$.
By~\eqref{eq:gap-small-scale}, the cubic term of $\Gamma_{E,T}(c,h)$
vanishes at this center, so the cubic expansion does not decide the sign.
We need an argument valid at every scale. Since $\Gamma_{E,T}(c,0)=0$, it is
enough to show that $\Gamma_{E,T}(c,\cdot)$ is strictly increasing on
$(0,\infty)$. Differentiation over the finite support gives
\begin{equation}\label{eq:binomial-scale-derivative}
 \frac{d}{dh}\Gamma_{E,T}(c,h)=\phi(0)D(h^2/2),
\end{equation}
where
\begin{equation}\label{eq:binomial-laplace}
 D(s)=\E[(Z-c)e^{-s(Z-c)^2}]+ce^{-sc^2}.
\end{equation}

Each support value enters $D$ through a signed weight $z-c$ and a squared
distance $(z-c)^2$. Accordingly, place an atom of weight
$\nu_z=\Pr\{Z=z\}(z-c)$ at $y_z=(z-c)^2$ for each support value $z$, and one
more atom of weight $c$ at $c^2$ for the second term. Then
$D(s)=\sum\nu e^{-sy}$ is the Laplace transform of the finite signed
measure~$\nu$. Its first two moments are the coefficients that vanished in the
expansion:
\begin{equation}\label{eq:binomial-cancellations}
 \sum\nu=\E Z=0,\qquad
 \sum\nu y=\E[(Z-c)^3]+c^3=\E Z^3-3c\E Z^2=0,
\end{equation}
so $D(0)=0$ because the count has mean $E$, and $D'(0)=0$ because of the
choice of $c$.

The positive and negative parts of $\nu$ therefore have equal mass and equal
first moment. $D(s)$ compares their integrals of the convex function
$y\mapsto e^{-sy}$, so the sign should come from which part is more spread
out. For measures of equal mass and first moment, spread is measured by the
hinge transform
\begin{equation}\label{eq:binomial-double-cumulative}
 J(t)=\sum\nu(t-y)_+,\qquad (u)_+=\max\{u,0\}:
\end{equation}
$J\geq0$ says that the positive part dominates the negative part in convex
order. We use $J$ only through an exact identity, which also yields
strictness. For $s>0$, finite summation and
\[
 e^{-sy}=s^2\int_y^\infty e^{-st}(t-y)\,dt
\]
give
\begin{equation}\label{eq:binomial-laplace-hinge}
 D(s)=s^2\int_0^\infty e^{-st}J(t)\,dt.
\end{equation}
By~\eqref{eq:binomial-cancellations}, $J(t)=t\sum\nu-\sum\nu y=0$ beyond
the largest location, so $J$ is continuous, piecewise affine and compactly
supported. The low-rate comparison thus reduces to showing that $J\geq0$
and $J\not\equiv0$.

To locate the kinks of $J$, note that the atoms at $c^2$ from $Z=0$ and from
the extra term combine to net weight $c(1-\Pr\{Z=0\})>0$. The other positive
weights come from displacements $z=j\geq1$ and sit at the \emph{positive
knots} $(j-c)^2$. The negative weights come from $z=-j\leq-1$ and sit at
$(j+c)^2$. Because $0<c<1/2$, the locations interlace:
\[
 c^2<(1-c)^2<(1+c)^2<(2-c)^2<(2+c)^2<\cdots.
\]
Between consecutive positive knots, $J$ has at most one kink, and its slope
there drops. So $J$ is concave on each such interval, and it suffices to check
$J\geq0$ at the positive knots. At a positive knot $t=(j-c)^2$, $J(t)$ is a
truncated cubic expectation over $|z-c|<j-c$. For $j\leq E$ we sum it
exactly. Put $C=T-E$, let $L\sim\operatorname{Binomial}(T-1,p)$, and write
$\ell_k=\Pr\{L=k\}$, with zero mass outside its support.
By~\eqref{eq:binomial-size-bias},
\[
 r(E+z)\Pr\{Z=z\}=Er\,\ell_{E+z-1},\qquad
 p(C-z)\Pr\{Z=z\}=Er\,\ell_{E+z}.
\]
Hence any summand of the form $r(E+z)A(z)-p(C-z)A(z+1)$, weighted by
$\Pr\{Z=z\}$, becomes $Er\{\ell_{E+z-1}A(z)-\ell_{E+z}A(z+1)\}$. Summed over
consecutive integers, it telescopes to two boundary terms. The expression
$r(E+z)A(z)-p(C-z)A(z+1)$ is the binomial Stein operator applied to $A$.
Our summand is cubic in $z$, so we look for a quadratic $A$.

\begin{lemma}[A binomial telescoping identity]\label{lem:binomial-stein-polynomial}
Put $C=T-E$. For $1\leq j\leq E$, define
\[
 R_j=j-c,\qquad A_j(z)=R_j^2-2Er-3c^2+z-z^2.
\]
Then, for every real $z$,
\begin{equation}\label{eq:binomial-cubic-stein-polynomial}
 \begin{split}
 &(z-c)\{R_j^2-(z-c)^2\}+c(R_j^2-c^2)\\
 &\hspace{2em}=r(E+z)A_j(z)-p(C-z)A_j(z+1).
 \end{split}
\end{equation}
\end{lemma}

\begin{proof}
Try $A(z)=\alpha+\beta z-z^2$ on the right and compare coefficients of the
two cubics, using $pT=E$, $rT=C$ and $Er=pC$. The $z^3$ coefficients are
both $-1$. The $z^2$ coefficient forces $\beta=2p+3c$, which equals $1$
precisely because $c=(1-2p)/3$. The $z$ coefficient then gives
$\alpha=R_j^2-2Er-3c^2$. Finally, the constant coefficients differ by
$Er(\beta-1)=0$. The critical center is thus used once, through $\beta=1$.
The identity holds for every real $R_j$; only the evaluation below uses
$R_j=j-c$.
\end{proof}

\begin{lemma}[Nonnegativity of the hinge transform]
\label{lem:binomial-double-cumulative}
If $T>2E$, then $J(t)\geq0$ for every $t\geq0$, and
$J((1-c)^2)>0$.
\end{lemma}

\begin{proof}
\emph{Positive knots with $j\leq E$.} Fix $1\leq j\leq E$ and
$t=R_j^2=(j-c)^2$. The atom at $y_z$ is active, $(t-y_z)_+>0$, exactly when
$|z-c|<j-c$, that is, $2c-j<z<j$; since $0<2c<1$, this means
$1-j\leq z\leq j-1$. The extra atom at $c^2$ is active as well. Splitting its
weight as $c=c\Pr\{|Z|\leq j-1\}+c\Pr\{|Z|\geq j\}$ puts the first part inside
the sum:
\[
 J((j-c)^2)=\sum_{z=1-j}^{j-1}\Pr\{Z=z\}
 \bigl[(z-c)\{R_j^2-(z-c)^2\}+c(R_j^2-c^2)\bigr]
 +c(R_j^2-c^2)\Pr\{|Z|\geq j\}.
\]
The bracket is the left side of~\eqref{eq:binomial-cubic-stein-polynomial}.
By Lemma~\ref{lem:binomial-stein-polynomial} and the size-bias identities,
the sum telescopes:
\begin{align*}
 &Er\sum_{z=1-j}^{j-1}
 \{\ell_{E+z-1}A_j(z)-\ell_{E+z}A_j(z+1)\}\\
 &\hspace{2em}=Er\{\ell_{E-j}A_j(1-j)-\ell_{E+j-1}A_j(j)\}.
\end{align*}
Both endpoint values equal
\[
 A_j(1-j)=A_j(j)=b_j:=j(1-2c)-2Er-2c^2,
\]
and $R_j^2-c^2=j(j-2c)$. Thus
\begin{equation}\label{eq:binomial-positive-knot}
 \begin{split}
 J((j-c)^2)={}&Er b_j(\ell_{E-j}-\ell_{E+j-1})\\
              &+cj(j-2c)\Pr\{|Z|\geq j\}.
 \end{split}
\end{equation}
The tail term is nonnegative. For the boundary term, $1-2c>0$ gives
$b_j\leq b_E$, and the identity $1-2r=-3c$ gives
\[
 b_E=-5Ec-2c^2<0.
\]
So the boundary term is nonnegative as soon as the lower mass of each pair
is the smaller one:
\begin{equation}\label{eq:binomial-pairing}
 \ell_{E-j}\leq\ell_{E+j-1}\qquad(1\leq j\leq E).
\end{equation}
At $j=1$ the two masses are equal. For $1\leq j<E$, the ratios taking the
upper and lower members of one pair to the next are
\[
 \frac{\ell_{E+j}}{\ell_{E+j-1}}
 =\frac{C-j}{E+j}\frac EC,\qquad
 \frac{\ell_{E-j-1}}{\ell_{E-j}}
 =\frac{E-j}{C+j}\frac CE.
\]
The first ratio is at least the second: after clearing positive
denominators, their difference has the sign of $j^2(C^2-E^2)\geq0$, since
$C>E$. Induction proves~\eqref{eq:binomial-pairing}, so
$J((j-c)^2)\geq0$ for $1\leq j\leq E$. At $j=1$ the boundary term vanishes,
while the tail term is $c(1-2c)\Pr\{|Z|\geq1\}>0$. Hence $J((1-c)^2)>0$.

\emph{Positive knots with $j>E$.} Now every negative displacement satisfies
$|z-c|\leq E+c<j-c$, so all atoms except the positive ones with $z\geq j$ are
active. Subtracting the vanishing full sum
$\sum\nu(t-y)=t\sum\nu-\sum\nu y=0$ leaves
\[
 J((j-c)^2)=\sum_{z\geq j}\Pr\{Z=z\}(z-c)
             \bigl((z-c)^2-(j-c)^2\bigr)\geq0.
\]

\emph{Between the knots.} Before $c^2$, $J$ vanishes. On
$[c^2,(1-c)^2]$ only the net atom at $c^2$ is active, so
\[
 J(t)=c(1-\Pr\{Z=0\})(t-c^2)\geq0
 \qquad(c^2\leq t\leq(1-c)^2).
\]
Between $(j-c)^2$ and $(j+1-c)^2$, $j\geq1$, the only possible interior knot
is $(j+c)^2$, with weight $(-j-c)\Pr\{Z=-j\}\leq0$. So $J$ is concave there
and lies above the chord joining its nonnegative endpoint values. Beyond the
largest location $J$ vanishes. This proves $J(t)\geq0$ for all $t\geq0$.
\end{proof}

By Lemma~\ref{lem:binomial-double-cumulative}, $J$ is continuous,
nonnegative, and positive at $(1-c)^2$. Hence the integral
in~\eqref{eq:binomial-laplace-hinge} is strictly positive for every $s>0$.
Equation~\eqref{eq:binomial-scale-derivative} and
$\Gamma_{E,T}(c,0)=0$ now give
\begin{equation}\label{eq:binomial-low-boundary-strict}
 \Gamma_{E,T}(c,h)>0\qquad(T>2E,\ h>0).
\end{equation}

\subsection*{Zero-center signs and reflection}
\begin{lemma}[The zero-center sign]\label{lem:binomial-reflected-zero}
For $h>0$, the gap $\Gamma_{E,T}(0,h)$ is negative when $T>2E$, zero when
$T=2E$, and positive when $E<T<2E$.
\end{lemma}

\begin{proof}
We first treat a general low-rate mean $m$, so that the same computation also
serves the high-rate case after reflection. Fix an integer $m$ with
$0<m<T/2$. Let $K\sim\operatorname{Binomial}(T,m/T)$ and
$L\sim\operatorname{Binomial}(T-1,m/T)$, and put
\[
 \psi_h(k)=\exp\left(-\frac{h^2(k-m)^2}{2}\right).
\]
The size-bias identities~\eqref{eq:binomial-size-bias} with $p=m/T$ give,
for every function $f$,
\[
 \E[(K-m)f(K)]
 =m(1-m/T)\E[f(L+1)-f(L)].
\]
Differentiating the gap and applying this identity gives
\begin{equation}\label{eq:binomial-reflected-stein}
 \frac{d}{dh}\Gamma_{m,T}(0,h)
 =\phi(0)m(1-m/T)\E[\psi_h(L+1)-\psi_h(L)].
\end{equation}
Pair $L=m-1-j$ with $L=m+j$ for $0\leq j<m$. Their increments are $d_j$
and $-d_j$, respectively, where
\[
 d_j=e^{-h^2j^2/2}-e^{-h^2(j+1)^2/2}>0.
\]
By~\eqref{eq:binomial-pairing}, applied with $E=m$, the upper mass is at
least the lower mass, so each pair contributes a nonpositive amount.
The remaining upper atoms, indexed by $m\leq j<T-m$, have positive mass
and negative increment $-d_j$. At least one exists because $m<T/2$.
Thus the derivative in~\eqref{eq:binomial-reflected-stein} is negative for
$h>0$. Since $\Gamma_{m,T}(0,0)=0$, the low-rate gap is negative.

If $T=2E$, the law of $Z$ is symmetric and Gaussian symmetry gives
$\Gamma_{E,T}(0,h)=0$. If $E<T<2E$, put $m=T-E<T/2$ and reflect the count:
$K'=T-K\sim\operatorname{Binomial}(T,m/T)$, with $K'-m=-Z$. Hence
\[
 \Gamma_{E,T}(0,h)=-\Gamma_{m,T}(0,h)>0.
\]
\end{proof}

\begin{proof}[Proof of Theorem~\ref{thm:rate-adaptive-center}]
For $T>2E$, equation~\eqref{eq:binomial-low-boundary-strict} proves positivity
at $\theta=c_{E,T}$, and Lemma~\ref{lem:binomial-propagation} extends it to
larger centers. For $E<T<2E$, Lemma~\ref{lem:binomial-reflected-zero} gives
positivity at zero, and the same propagation lemma handles every positive
center. At $T=2E$, symmetry gives equality at zero and propagation gives
strict positivity for $\theta>0$.
\end{proof}

\subsection*{Strict comparison at the symmetric endpoint}
At the symmetric endpoint of Theorem~\ref{thm:rate-adaptive-pointwise}, the
full-batch threshold gives a tie. A slightly larger threshold already
gives a strict improvement using the average alone.

\begin{lemma}[Improving the threshold at half rate]\label{lem:strict-average-endpoint}
Let $E\in\N_{>0}$, $a>0$, and $K\sim\operatorname{Binomial}(2E,1/2)$.
Let $Q=\Normal(0,1)$, and let $P$ be the mixture whose conditional law
given $K=k$ is $\Normal(ak/E,1)$. With $\varepsilon=a^2/2$,
\[
 H_{e^\varepsilon}(P, Q)
 >\frac12-e^\varepsilon\overline\Phi(a).
\]
\end{lemma}

\begin{proof}
Symmetry of $K-E$ and of the Gaussian noise makes $P$ symmetric about
$a$, with no atom there. Thus $P((a,\infty))=1/2$, and the value of this
ray is $\delta_*:=1/2-e^\varepsilon\overline\Phi(a)$.
The likelihood ratio of the mixture is the continuous function
\[
 L(y)=\E\exp\!\left(\frac{aK}{E}y-\frac{a^2K^2}{2E^2}\right).
\]
Completing the square at the threshold gives
\[
 L(a)=e^{a^2/2}\E\exp\!\left[
       -\frac{a^2}{2}(K/E-1)^2\right]<e^{a^2/2}=e^\varepsilon.
\]
The inequality is strict because $a>0$ and $\Pr\{K\ne E\}>0$.
By continuity there is $\eta>0$ such that $L(y)<e^\varepsilon$ for
$a\leq y\leq a+\eta$. Since $\phi(y)>0$,
\[
 P((a+\eta,\infty))-e^\varepsilon Q((a+\eta,\infty))
 =\delta_*-\int_a^{a+\eta}(L(y)-e^\varepsilon)\phi(y)\,dy
 >\delta_*.
\]
Taking the supremum over events proves the claim.
\end{proof}

\subsection*{Sharpness uniformly over the Gaussian scale}
Equation~\eqref{eq:gap-small-scale} identifies the cubic coefficient.
Since $\operatorname{Var}(Z)>0$, every $\theta<c_{E,T}$ gives a negative gap
for sufficiently small $h>0$. Together with
Theorem~\ref{thm:rate-adaptive-center}, this proves that, in the low-rate
regime, $\theta\geq c_{E,T}$ is necessary and sufficient for positivity at
every $h>0$. This is sharpness of the scale-uniform sum-event comparison,
not a classification of the full transcript profile when the scale and
privacy parameters are coupled.

\section{The fixed-participation privacy limit}\label{app:sparse-limit}
We first prove the profile limit for $Tq_T\to\rho>0$, then invert it to
obtain the calibrated-noise offset.

\begin{theorem}[The fixed-participation profile limit]\label{thm:accountant-limit}
Let $\rho>0$, $\Delta>0$, and $\varepsilon\geq0$. Suppose
$q_T\in[0,1]$, $Tq_T\to\rho$, $\sigma_T>0$ eventually, and
\[
 \frac{\Delta}{\sigma_T}-\sqrt{2\log T}\longrightarrow s\in\R.
\]
Then
\[
 \AO_S(\sigma_T,\Delta,q_T,T,\varepsilon)
 \longrightarrow1-\exp\{-\rho\Phi(s)\}.
\]
\end{theorem}

\begin{proof}
Put $\mu_T=\Delta/\sigma_T$ and $b_T=\sqrt{2\log T}$. Standardizing every
coordinate by $\sigma_T$ makes the reference observations independent
$Z_i\sim\Normal(0,1)$. Their likelihood ratio is
\begin{equation}\label{eq:sparse-likelihood-ratio}
 R_T=\prod_{i=1}^T
 \left[1-q_T+q_T\exp\left(\mu_TZ_i-\frac{\mu_T^2}{2}\right)\right].
\end{equation}

Write $Y_T(z)=\exp(\mu_Tz-\mu_T^2/2)$. We linearize $\log R_T$ after
truncating each coordinate at a level $t_T$.
The level has two jobs. It must be high enough that, with reference
probability tending to one, no coordinate exceeds it, so truncation changes
nothing. We therefore need $T\overline\Phi(t_T)\to0$; Mills' bound makes
$T/t_T=o(e^{t_T^2/2})$ sufficient. The level must also be low enough that every
perturbation $q_T(Y_T-1)$ of the likelihood ratio is uniformly small and the
perturbations have vanishing total second moment. Below this level,
$Y_T\leq e^{t_T^2/2}$; together with the moment bound below, this makes
$e^{t_T^2/2}=o(T)$ sufficient. The naive level $t_T=b_T$ gives
$e^{t_T^2/2}=T$, and our coarse second-moment bound stays of order one.
Lowering the level by $\log b_T/(2b_T)$ gives
$e^{t_T^2/2}=Tb_T^{-1/2}(1+o(1))$, which lies between $T/t_T$ and $T$. The first
truncated moment then produces the limit, and the second moment controls both
the fluctuation of the linear term and the linearization error.

\paragraph{The reference maximum.}
For all sufficiently large $T$, define
\[
 t_T=b_T-\frac{\log b_T}{2b_T}.
\]
Mills' upper bound gives
\[
 T\Pr\{Z>t_T\}\leq T\frac{\phi(t_T)}{t_T}\longrightarrow0,
\]
because
\[
 \log T-\frac{t_T^2}{2}
 =\frac{\log b_T}{2}-\frac{(\log b_T)^2}{8b_T^2}.
\]
The upper bound is of order $b_T^{-1/2}$. Thus every coordinate lies below
$t_T$ with probability tending to one. We will also use
\begin{equation}\label{eq:truncation-scales}
 \frac{e^{t_T^2/2}}{T}=b_T^{-1/2}(1+o(1)),\qquad
 q_Te^{t_T^2/2}\to0,\qquad Tq_T^2e^{t_T^2/2}\to0.
\end{equation}
The last two estimates follow by multiplying the first by $Tq_T\to\rho$
and by $(Tq_T)^2$, respectively.

\paragraph{Truncated moments.}
Define
\[
 X_T(z)=q_T(Y_T(z)-1)\mathbf1_{\{z\leq t_T\}}.
\]
Completing the square gives
\begin{align}
 \E[Y_T(Z)\mathbf1_{\{Z\leq t_T\}}]
 &=\Phi(t_T-\mu_T),\label{eq:tilted-first}\\
 \E[Y_T(Z)^2\mathbf1_{\{Z\leq t_T\}}]
 &=e^{\mu_T^2}\Phi(t_T-2\mu_T).\label{eq:tilted-second}
\end{align}
Hence
\[
 T\E X_T(Z)=Tq_T[\Phi(t_T-\mu_T)-\Phi(t_T)]
 \longrightarrow-\rho\Phi(s),
\]
since $t_T-\mu_T\to-s$ and $\Phi(t_T)\to1$.

Because $Y_T\geq0$, $(Y_T-1)^2\leq Y_T^2+1$, so
\begin{equation}\label{eq:truncated-second-bound}
 T\E X_T(Z)^2\leq Tq_T^2
 \left[e^{\mu_T^2}\Phi(t_T-2\mu_T)+\Phi(t_T)\right].
\end{equation}
The second term tends to zero. For the first,
$2\mu_T-t_T=b_T+2s+o(1)\to\infty$, and eventually $2\mu_T-t_T\geq1$.
Mills' bound and the identity
\[
 \mu_T^2-\frac{(2\mu_T-t_T)^2}{2}
 =\frac{t_T^2}{2}-(\mu_T-t_T)^2
\]
therefore imply
\[
 e^{\mu_T^2}\Phi(t_T-2\mu_T)
 \leq\frac{e^{t_T^2/2}}{\sqrt{2\pi}}.
\]
By~\eqref{eq:truncation-scales}, the right side
of~\eqref{eq:truncated-second-bound} tends to zero. Independence gives
\[
 \operatorname{Var}\left(\sum_{i=1}^T X_T(Z_i)\right)
 \leq T\E X_T(Z)^2\longrightarrow0.
\]
Chebyshev's inequality yields
\begin{equation}\label{eq:linear-sum-limit}
 \sum_{i=1}^T X_T(Z_i)\longrightarrow-\rho\Phi(s)
 \quad\text{in reference probability}.
\end{equation}

\paragraph{The log likelihood ratio.}
Let $A_T=\{\max_i Z_i\leq t_T\}$, so $\Pr(A_T^c)\to0$.
Eventually $\mu_T\geq0$, and for $z\leq t_T$,
\[
 |q_T(Y_T(z)-1)|
 \leq q_T+q_T\exp\left(\mu_Tt_T-\frac{\mu_T^2}{2}\right)
 \leq q_T+q_Te^{t_T^2/2}\longrightarrow0.
\]
The bound is uniform in $z\leq t_T$. Eventually every truncated perturbation
therefore has absolute value at most $1/2$, where
$|\log(1+x)-x|\leq x^2$. On $A_T$,
\begin{align*}
 \mathbf1_{A_T}\left|\log R_T-\sum_{i=1}^T X_T(Z_i)\right|
 &\leq\mathbf1_{A_T}\sum_{i=1}^T
   |\log(1+X_T(Z_i))-X_T(Z_i)|\\
 &\leq\sum_{i=1}^T X_T(Z_i)^2.
\end{align*}
The expectation of the last expression tends to zero by
\eqref{eq:truncated-second-bound}. Markov's inequality, the bound on
$\Pr(A_T^c)$, and~\eqref{eq:linear-sum-limit} imply
\begin{equation}\label{eq:log-likelihood-limit}
 \log R_T\longrightarrow-\rho\Phi(s)
 \quad\text{in reference probability}.
\end{equation}
Thus $R_T\to c_s:=e^{-\rho\Phi(s)}$ in reference probability.
Although $\E R_T=1$ for every $T$, $c_s<1$: rare large values retain the
missing mean, so convergence of unbounded expectations cannot be used.

\paragraph{The directed privacy limits.}
Let $\alpha=e^\varepsilon\geq1$. Writing $P_T,Q_T$ for the two transcript
laws, the directed profiles are
\begin{align*}
 H_\alpha(P_T, Q_T)&=1-\E_{Q_T}\min\{R_T,\alpha\},\\
 H_\alpha(Q_T, P_T)&=1-\alpha\E_{Q_T}\min\{R_T,\alpha^{-1}\}.
\end{align*}
Convergence in probability to a constant implies convergence of expectations
for bounded continuous functions. Applying this fact to the two truncations
gives
\[
 H_\alpha(P_T, Q_T)\to1-c_s,\qquad
 H_\alpha(Q_T, P_T)\to(1-\alpha c_s)_+.
\]
Since $(1-\alpha c_s)_+\leq1-c_s$, taking the maximum proves the theorem.
\end{proof}

\begin{theorem}[Fixed-participation calibration]
\label{thm:fixed-participation-calibration}
Let $\rho>0$, $\Delta>0$, $\varepsilon\geq0$, and
$0<\delta<1-e^{-\rho}$. Suppose $q_T\in[0,1]$ and $Tq_T\to\rho$, and define
\[
 \sigma_T^{\mathrm{cal}}:=\inf\left\{\sigma\geq0:
 \AO_S(\sigma,\Delta,q_T,T,\varepsilon)\leq\delta\right\}.
\]
Then $\sigma_T^{\mathrm{cal}}>0$ eventually, and
\[
 \frac{\Delta}{\sigma_T^{\mathrm{cal}}}-\sqrt{2\log T}
 \longrightarrow s_{\rho,\delta}:=
 \Phi^{-1}\!\left(-\frac{\log(1-\delta)}\rho\right).
\]
For the capped exact-budget schedule $q_T=\min\{1,\rho/T\}$, it follows that
\[
 \frac{\sigma_T^{\mathrm{cal}}\sqrt{2\log(1/q_T)}}\Delta\longrightarrow1.
\]
\end{theorem}

\begin{proof}
The function $L_\rho(s)=1-e^{-\rho\Phi(s)}$ is continuous and strictly
increasing, with range $(0,1-e^{-\rho})$ and inverse value
$L_\rho(s_{\rho,\delta})=\delta$. For a fixed offset $u$, consider the
eventually positive candidate
\[
 \sigma_T(u)=\frac\Delta{b_T+u},\qquad b_T=\sqrt{2\log T}.
\]
Given $\eta>0$,
\[
 L_\rho(s_{\rho,\delta}-\eta)<\delta
 <L_\rho(s_{\rho,\delta}+\eta).
\]
Theorem~\ref{thm:accountant-limit} makes the higher-noise candidate
$\sigma_T(s_{\rho,\delta}-\eta)$ eventually feasible and the lower-noise
candidate $\sigma_T(s_{\rho,\delta}+\eta)$ eventually infeasible.
The finite calibration properties proved in Appendix~\ref{app:calibration}
apply once $q_T>0$, which holds eventually. Monotonicity and feasibility give
\[
 \frac\Delta{b_T+s_{\rho,\delta}+\eta}
 \leq\sigma_T^{\mathrm{cal}}
 \leq\frac\Delta{b_T+s_{\rho,\delta}-\eta}.
\]
In particular, the calibrated noise is eventually positive. Taking
reciprocals gives
\[
 s_{\rho,\delta}-\eta
 \leq\frac\Delta{\sigma_T^{\mathrm{cal}}}-b_T
 \leq s_{\rho,\delta}+\eta.
\]
Since $\eta$ is arbitrary, the offset limit follows. For the capped schedule,
$q_T=\rho/T$ eventually and $\log(1/q_T)/\log T\to1$, giving the final claim.
\end{proof}

Equivalently, the calibrated standard deviation satisfies
\begin{equation}\label{eq:offset-approximation}
 \sigma_T^{\mathrm{cal}}
 =\frac\Delta{\sqrt{2\log T}+s_{\rho,\delta}+o(1)}.
\end{equation}

\begin{proof}[Proof of Theorem~\ref{thm:smallbatch}]
Apply Theorem~\ref{thm:fixed-participation-calibration} with $\rho=E$.
For $T\geq E$, its capped schedule is exactly $q=E/T$, so the offset limit
and~\eqref{eq:sigmaasymptotic} follow. Squaring and multiplying by $T/E^2$
gives~\eqref{eq:privacydivergence}. The sampling term in~\eqref{eq:variance}
converges to $S/E$, while the privacy variance diverges. Their sum is
therefore asymptotic to $v_{E,T}$, proving~\eqref{eq:total-equivalent}.
\end{proof}

\section{Rational bounds for the finite reversal}\label{app:rational-calculation}
This appendix proves the privacy bounds used in
Theorem~\ref{thm:strict-interior-counterexample}. We partition the first two
reference coordinates into rectangles and integrate over the third
coordinate exactly. Rational interval arithmetic then bounds both directed
profiles. The supplement~\cite{rational-supplement} contains the individual
logarithm brackets and unrounded interval expressions. Decimal values, where
given, are for orientation only.

\subsection{Conditional tails and rectangle bounds}
Fix $E=2$, $T=3$, $\Delta=1$, and $\varepsilon=1/4$. The candidate
privacy variance $v=5/32$ corresponds to
\[
 \sigma^2=\frac{E^2v}{T}=\frac5{24},\qquad
 \mu=\frac\Delta\sigma=\sqrt{\frac{24}{5}}.
\]
Write $P=P_{2/3,\sqrt{5/24}}$, $Q=Q_{\sqrt{5/24}}$, $q=2/3$, and
$\alpha=e^{1/4}$. After standardization, the one-coordinate likelihood ratio
under the reference law is
\[
 r(z)=\frac13+\frac23\exp\left(\mu z-\frac{\mu^2}{2}\right).
\]
Condition on the first two reference coordinates and put
$A=r(Z_1)r(Z_2)$. The remaining conditional hinges are
\[
 F(A)=\E[(Ar(Z)-\alpha)_+],\qquad
 G(A)=\E[(1-\alpha Ar(Z))_+],
\]
where $Z$ is standard normal. If $\alpha>A(1-q)$, define
\[
 c_F(A)=\frac{\log((\alpha/A-(1-q))/q)+\mu^2/2}{\mu}.
\]
Gaussian tail integration gives
\[
 F(A)=\bigl(A(1-q)-\alpha\bigr)\overline\Phi(c_F(A))
       +Aq\,\overline\Phi(c_F(A)-\mu).
\]
If $\alpha\leq A(1-q)$, then $F(A)=A-\alpha$. Similarly, if
$1>\alpha A(1-q)$, define
\[
 c_R(A)=\frac{\log((1/(\alpha A)-(1-q))/q)+\mu^2/2}{\mu}.
\]
Then
\[
 G(A)=\bigl(1-\alpha A(1-q)\bigr)\Phi(c_R(A))
       -\alpha Aq\,\Phi(c_R(A)-\mu),
\]
whereas $G(A)=0$ if $1\leq\alpha A(1-q)$.

Both $F$ and $G$ are convex in $A$, being expectations of positive parts of
affine functions. Let $\gamma=\Normal(0,1)$. On a rectangle $C$ where
$L\leq A\leq U$ with $L<U$, convexity gives
\begin{equation}\label{eq:rectangle-chord}
 \int_C F(A)\,d\gamma^{\otimes2}
 \leq\frac{UQ_C-P_C}{U-L}F(L)
      +\frac{P_C-LQ_C}{U-L}F(U),
\end{equation}
where $Q_C=\gamma^{\otimes2}(C)$ and
$P_C=\int_C A\,d\gamma^{\otimes2}$. The same bound holds for $G$.
For $C=(x_i,x_{i+1}]\times(x_j,x_{j+1}]$, put
\begin{align*}
 Q_i&=\Phi(x_{i+1})-\Phi(x_i),\\
 P_i&=(1-q)Q_i+q[\Phi(x_{i+1}-\mu)-\Phi(x_i-\mu)].
\end{align*}
Then $Q_C=Q_iQ_j$, $P_C=P_iP_j$, and monotonicity of $r$ gives
\[
 L=r(x_i)r(x_j),\qquad U=r(x_{i+1})r(x_{j+1}).
\]
We insert outward bounds on these quantities into~\eqref{eq:rectangle-chord}.
Outside the rectangle union, the bounds $F(A)\leq A$ and $G(A)\leq1$ control
the complement contributions.

The eleven rational edges are
\[
 -\frac{31}{10},\ \frac35,\ \frac{21}{20},\ \frac{27}{20},\
 \frac{33}{20},\ \frac{19}{10},\ \frac{43}{20},\ \frac{12}{5},\
 \frac{53}{20},\ \frac{57}{20},\ \frac{31}{10}.
\]
They form one hundred rectangles. The next subsection specifies the
interval rules; Table~\ref{tab:rational-cells} gives the resulting bounds.

\subsection{Elementary interval rules}
All interval endpoints are rational. Addition and subtraction use the usual
outward endpoint rules; multiplication takes the minimum and maximum of the
four endpoint products; reciprocation of $[l,u]$ with $l>0$ gives
$[1/u,1/l]$. The required signs are checked before reciprocation.

\paragraph{Exponential and logarithm.}
Choose integers $b\geq0$, $n\geq1$ with $|x|\leq2^b$, and put
$y=|x|/2^b$. Define
\[
 L_n(y)=\sum_{k=0}^{n-1}\frac{y^k}{k!},\qquad
 U_n(y)=L_n(y)+\frac{y^n(n+1)}{n!\,n}.
\]
For $0\leq y\leq1$,
$L_n(y)\leq e^y\leq U_n(y)$: the tail starting at order $n$ is bounded by
$y^n/n!$ times a geometric sum with ratio $1/(n+1)$. Raise the endpoints to
the power $2^b$ and, for negative $x$, reciprocate with reversed endpoints.
Monotonicity handles interval inputs. The calculation uses $(b,n)=(1,20)$
for $e^{1/4}$, $(5,20)$ for the likelihood-ratio edges, and $(6,24)$ for
Gaussian density exponentials.

A logarithm bracket $[l,u]$ for a positive interval $[a,b]$ is accepted only
if the exponential upper bound at $l$ is at most $a$ and the exponential
lower bound at $u$ is at least $b$. Monotonicity of $\log$ then proves the
enclosure. The supplement lists the individual rational brackets and Taylor
orders; floating-point logarithms do not justify any bound.

\paragraph{Square roots.}
For an integer $m>0$, let $N=10^6$ and choose the integer $j$ satisfying
$j^2\leq mN^2<(j+1)^2$. Then
$\sqrt m\in[j/N,(j+1)/N]$. Reciprocal interval arithmetic gives the required
quantities
\[
 \mu=\frac{24}{\sqrt{120}},\qquad
 \mu^{-1}=\frac5{\sqrt{120}}.
\]

\paragraph{Central Gaussian CDF values.}
Use the enclosure
\[
 n_-:=\frac{39894228040143}{10^{14}}
 <\phi(0)<
 n_+:=\frac{39894228040144}{10^{14}}.
\]
For $0\leq x\leq7/5$, put
\[
 P_0(x)=x-\frac{x^3}{6}+\frac{x^5}{40}-\frac{x^7}{336}
       +\frac{x^9}{3456}-\frac{x^{11}}{42240}+\frac{x^{13}}{599040}.
\]
Taylor's theorem for $e^{-u}$, followed by integration, gives
\[
 \left|\frac{\Phi(x)-1/2}{\phi(0)}-P_0(x)\right|
 \leq\frac{x^{15}}{2^7 7!\,15}
 \leq\frac{x^{15}}{8467200}.
\]
Here $2^7 7!\,15=9676800$; the final denominator deliberately gives a
weaker rational bound. Interval multiplication with $[n_-,n_+]$ gives the central CDF bounds.
The supplement includes the normalizer enclosure and its rational square
comparisons.

\paragraph{Gaussian tails.}
For $x>0$, define
\[
 d_{k,0}(x)=x,\qquad
 d_{k,n+1}(x)=x+\frac{k}{d_{k+1,n}(x)},\qquad
 M_n(x)=\frac1{d_{1,n}(x)}.
\]
The Mills bounds used in the calculation are
\[
 \phi(x)M_{13}(x)\leq\Phi(-x)\leq\phi(x)M_{14}(x).
\]
To verify their directions, let
\[
 I_k(x)=\int_0^\infty t^k e^{-xt-t^2/2}\,dt.
\]
Integration by parts gives $1=xI_0+I_1$ and
$kI_{k-1}=xI_k+I_{k+1}$ for $k\geq1$. Every ratio $I_k/I_{k-1}$ is positive.
Iterating these identities expresses $I_0=\Phi(-x)/\phi(x)$ as the continued
fraction above with a positive last remainder. Replacing that remainder by
zero alternates the direction at each reciprocal, giving a lower bound at
odd depth and an upper bound at even depth.

The calculation uses these tail bounds for $x>7/5$. Negative arguments are
handled by $\Phi(-x)=1-\Phi(x)$, and interval arguments by evaluating the
appropriate endpoints.

\subsection{The certificate and final bounds}
Let $x_0,\ldots,x_{10}$ be the eleven edges above in increasing order.
For each rectangle, the table gives upper-bound numerators with denominator
$10^{10}$, obtained by rounding the rational chord bound upward. The bounds
are symmetric in $i,j$, so only $i\leq j$ is shown; each off-diagonal row
is counted twice.

The outward complement bounds are
\[
 B_+=\frac{2286836197}{10^{10}},\qquad
 B_-=\frac{38961008}{10^{10}}.
\]
They round upward the rational bounds
\[
 \frac{223323847325382457311}{976562500000000000000},\qquad
 \frac{974025185261562567831}{250000000000000000000000},
\]
respectively. Adding the table entries, with off-diagonal multiplicity two,
and the complement contributions yields
\begin{align}
 H_{e^{1/4}}(P^{\otimes3}, Q^{\otimes3})
 &\leq\frac{7420364678}{10^{10}}<\frac34,
 \label{eq:rational-forward-total}\\
 H_{e^{1/4}}(Q^{\otimes3}, P^{\otimes3})
 &\leq\frac{7237647581}{10^{10}}<\frac34.
 \label{eq:rational-reverse-total}
\end{align}
Thus the sampled candidate is feasible, and $v_3\leq5/32$.

For the full-batch lower bracket, evaluate the profile at $v=4/25$. Then
$a=5/2$, and the Gaussian formula gives
\[
 G_{1/4}(5/2)=\Phi(23/20)-e^{1/4}\Phi(-27/20).
\]
The central CDF bounds and the exponential interval with $(b,n)=(1,20)$ give
\[
 \Phi(23/20)-e^{1/4}\Phi(-27/20)>\frac{19}{25}>\frac34.
\]
Hence $4/25$ is infeasible at full batch, and
\eqref{eq:calibration-transfer} gives $4/25<v_2$. Together with the sampled
bounds, this proves the privacy-variance inequalities in
Theorem~\ref{thm:strict-interior-counterexample}.

\begingroup\small
\begin{longtable}{@{}rrrr@{}}
\caption{Directed rectangle bounds, with denominator $10^{10}$.
Off-diagonal rows are counted twice.}
\label{tab:rational-cells}\\
\toprule $i$ & $j$ & Forward numerator & Reverse numerator \\\midrule
\endfirsthead
\toprule $i$ & $j$ & Forward numerator & Reverse numerator \\\midrule
\endhead
\midrule
\multicolumn{4}{r}{\small Continued on the next page}\\
\endfoot
\bottomrule
\endlastfoot
0 & 0 & 179309587 & 4541683687 \\
0 & 1 & 75622963 & 718523085 \\
0 & 2 & 69056184 & 283059965 \\
0 & 3 & 91045268 & 147341455 \\
0 & 4 & 94610823 & 49317043 \\
0 & 5 & 113839529 & 15318582 \\
0 & 6 & 131941862 & 0 \\
0 & 7 & 139954908 & 0 \\
0 & 8 & 111297082 & 0 \\
0 & 9 & 125975330 & 0 \\
1 & 1 & 30525041 & 104528577 \\
1 & 2 & 27152594 & 35842913 \\
1 & 3 & 35097753 & 15069211 \\
1 & 4 & 36770950 & 3637515 \\
1 & 5 & 44506072 & 0 \\
1 & 6 & 49279269 & 0 \\
1 & 7 & 50820946 & 0 \\
1 & 8 & 38394831 & 0 \\
1 & 9 & 42280938 & 0 \\
2 & 2 & 23806351 & 9614344 \\
2 & 3 & 31009020 & 3290409 \\
2 & 4 & 32740303 & 0 \\
2 & 5 & 37920442 & 0 \\
2 & 6 & 41171906 & 0 \\
2 & 7 & 40662238 & 0 \\
2 & 8 & 30071802 & 0 \\
2 & 9 & 32717715 & 0 \\
3 & 3 & 40879835 & 59609 \\
3 & 4 & 40928790 & 0 \\
3 & 5 & 46327961 & 0 \\
3 & 6 & 48166599 & 0 \\
3 & 7 & 46477322 & 0 \\
3 & 8 & 33960816 & 0 \\
3 & 9 & 36690918 & 0 \\
4 & 4 & 40398228 & 0 \\
4 & 5 & 43807871 & 0 \\
4 & 6 & 44478157 & 0 \\
4 & 7 & 42346729 & 0 \\
4 & 8 & 30722365 & 0 \\
4 & 9 & 33052386 & 0 \\
5 & 5 & 46472095 & 0 \\
5 & 6 & 46582913 & 0 \\
5 & 7 & 44021815 & 0 \\
5 & 8 & 31809211 & 0 \\
5 & 9 & 34139592 & 0 \\
6 & 6 & 46337541 & 0 \\
6 & 7 & 43592151 & 0 \\
6 & 8 & 31420920 & 0 \\
6 & 9 & 33672920 & 0 \\
7 & 7 & 40898832 & 0 \\
7 & 8 & 29435883 & 0 \\
7 & 9 & 31517496 & 0 \\
8 & 8 & 21168396 & 0 \\
8 & 9 & 22654198 & 0 \\
9 & 9 & 24237093 & 0 \\
\end{longtable}
\endgroup

\section{Elementary sufficient conditions}\label{app:elementary}
The next two arguments give simpler sufficient conditions for finite-horizon
dominance. Both use the sum event from the proof of
Theorem~\ref{thm:rate-adaptive-pointwise} and are subsumed by that theorem.

\begin{proposition}[Full-support convexity]\label{thm:near-full-accountant}
Let $E,T\in\N_{>0}$ with $T\geq E$, $\Delta>0$, $\varepsilon\geq0$, and
$0<\delta<1$. If
\begin{equation}\label{eq:nearfull-region}
 \frac TE\leq\frac12+\frac{\varepsilon v_{\mathrm{FB}}}{\Delta^2},
\end{equation}
then $\mathcal A_{E,T}(v_{\mathrm{FB}};\varepsilon)\geq\delta$.
For $T>E$, the inequality is strict, and
\[
 v_{E,T}>v_{\mathrm{FB}},\qquad
 V_E(E,S)<V_E(T,S)\quad(S\geq0).
\]
\end{proposition}

\begin{proof}
Put $a=\Delta/\sqrt{v_{\mathrm{FB}}}$ and $c_*=a/2+\varepsilon/a$.
At variance $v_{\mathrm{FB}}$, the sum event has sampled value
\[
 \E\left[\overline\Phi\left(c_*-\frac{aK}{E}\right)\right]
 -e^\varepsilon\overline\Phi(c_*),
 \qquad K\sim\operatorname{Binomial}(T,E/T).
\]
Define $f(x)=\overline\Phi(c_*-ax)$. Since
\[
 f''(x)=a^2(c_*-ax)\phi(c_*-ax),
\]
condition~\eqref{eq:nearfull-region} makes $f$ convex on $[0,T/E]$.
Jensen's inequality and $\E[K/E]=1$ therefore bound the event value below
by~\eqref{eq:fullbatch-event-identity}.

If $T>E$, the count is nondegenerate. Even at equality
in~\eqref{eq:nearfull-region}, $f''$ is positive before the right endpoint,
so $f$ is strictly convex on the whole interval. Strict Jensen makes the
sampled profile exceed $\delta$. Equation~\eqref{eq:calibration-transfer}
and the nonnegative sampling term give the variance comparisons.
\end{proof}

\begin{theorem}[Finite-support secants]\label{thm:support-width-accountant}
Let $T>E>0$ be integers, $\Delta>0$, $\varepsilon>0$, and $0<\delta<1$.
If
\begin{equation}\label{eq:support-width-region}
 a^2:=\frac{\Delta^2}{v_{\mathrm{FB}}}
 \leq\frac{2\varepsilon E}{T-1},
\end{equation}
then $\mathcal A_{E,T}(v_{\mathrm{FB}};\varepsilon)>\delta$, and
\[
 v_{E,T}>v_{\mathrm{FB}},\qquad
 V_E(E,S)<V_E(T,S)\quad(S\geq0).
\]
\end{theorem}

\begin{proof}
Put
\[
 u=\frac a2-\frac\varepsilon a,\qquad h=\frac aE,\qquad
 g(k)=\Phi(u+h(k-E)),\quad k=0,\ldots,T.
\]
As in~\eqref{eq:event-profile-bound}, it is enough to prove
\begin{equation}\label{eq:support-width-expectation}
 \E[g(K)]>g(E),\qquad K\sim\operatorname{Binomial}(T,E/T),
\end{equation}
except for an endpoint where changing the ray's threshold supplies strictness.

First suppose~\eqref{eq:support-width-region} is strict. Since $T-1\geq E$,
we have $a^2<2\varepsilon$ and $u<0$. Let $\ell$ be the affine line through
$g(E-1)$ and $g(E)$:
\[
 \ell(k)=g(E)+(k-E)[g(E)-g(E-1)].
\]
Convexity on the negative half-line gives $g(k)\geq\ell(k)$ for every
integer $0\leq k\leq E$.

For $k=E+d$, with $1\leq d\leq T-E$, and $0<t<1$,
\begin{align*}
 (u-ht)^2-(u+hdt)^2
 &=ht(d+1)[-2u-h(d-1)t],\\
 aE[-2u-h(d-1)t]
 &=2\varepsilon E-a^2[E+(d-1)t]\\
 &\geq2\varepsilon E-a^2(T-1)>0.
\end{align*}
Thus $\phi(u+hdt)>\phi(u-ht)$. Integrating yields
\begin{align*}
 g(E+d)-g(E)
 &=dh\int_0^1\phi(u+hdt)\,dt\\
 &>dh\int_0^1\phi(u-ht)\,dt
 =d[g(E)-g(E-1)].
\end{align*}
Hence $g(k)\geq\ell(k)$ on the entire support, with strict inequality at
$k=T$. Since $\E K=E$ and $\Pr\{K=T\}>0$,
\[
 \E[g(K)]>\E[\ell(K)]=g(E).
\]

At equality in~\eqref{eq:support-width-region}, put $m=T-E$. Then
$u=-h(m-1)/2$, and
\[
 -2u-h(d-1)t=h[(m-1)-(d-1)t].
\]
For $m>1$, this is positive whenever $0<t<1$ and $1\leq d\leq m$, so the
same argument applies. For $m=1$, $u=0$ and the right adjacent secant ties
by symmetry. Every support value still lies on or above $\ell$. If $E\geq2$,
strict convexity on the negative half-line gives
$g(E-2)>\ell(E-2)$, and that count has positive probability.
Thus~\eqref{eq:support-width-expectation} remains strict.

The remaining case is $E=1,T=2$, where $a^2=2\varepsilon$ and the selected
ray ties. Lemma~\ref{lem:strict-average-endpoint} improves its threshold,
so even the average-only forward profile is strictly larger than $\delta$.
In every case, equation~\eqref{eq:calibration-transfer} yields
$v_{E,T}>v_{\mathrm{FB}}$, and adding the sampling term gives the final claim.
\end{proof}

\paragraph{Comparison of the sufficient conditions.}
The coefficients multiplying $\varepsilon$ in these two conditions satisfy
\[
 B_{\mathrm{conv}}:=\frac{2E}{2T-E}
 <B_{\mathrm{sec}}:=\frac{2E}{T-1}
 \leq\Lambda_{E,T}\qquad(T>E).
\]
The first inequality follows from $2T-E>T-1$. If $E<T\leq2E$, then
$B_{\mathrm{sec}}\leq2=\Lambda_{E,T}$, with equality exactly at $T=E+1$.
If $T>2E$, integrality gives $T-1\geq2E$, so
\[
 B_{\mathrm{sec}}\leq1<\frac{6E}{3E+2}<\Lambda_{E,T}.
\]
Thus the secant condition improves on full-support convexity, but neither
condition enlarges the coverage of Theorem~\ref{thm:rate-adaptive-pointwise}.

In terms of the sampling rate, the secant condition is
\[
 q\geq\frac{a^2E}{a^2+2\varepsilon E}.
\]
Strict monotonicity of $G_\varepsilon$ gives the equivalent target condition
\[
 \delta\leq
 G_\varepsilon\!\left(\sqrt{\frac{2\varepsilon E}{T-1}}\right).
\]
For the nearest competitor $T=E+1$, it reduces to $a^2\leq2\varepsilon$, or
\[
 \delta\leq\frac12-e^\varepsilon\Phi(-\sqrt{2\varepsilon}).
\]
At $\varepsilon=0$ the secant condition has no solutions. These are
sufficient conditions; failure of a condition does not imply that a smaller
batch has lower variance.

\section{A forward integer-order R\'enyi comparison}\label{app:renyi}
A global comparison is available for forward integer-order R\'enyi
divergence~\cite{mironov2017}. Fix a privacy variance $v>0$. The raw variance
at $(E,T)$ is $Eqv$. For an integer $m\geq2$, put
$z=\Delta^2/(Eqv)$. The one-step forward likelihood moment is
\begin{equation}\label{eq:moment}
 M_m(q,z)=\sum_{k=0}^m\binom{m}{k}q^k(1-q)^{m-k}
          \exp\left(\frac{k(k-1)z}{2}\right),
\end{equation}
and the transcript moment is $M_m(q,z)^T$. Exact sampled-Gaussian moments
of this form are studied by Mironov, Talwar, and
Zhang~\cite{mironovtalwarzhang2019}. More generally, for $P\ll Q$, write
\[
 \mathcal P_m(P\mid Q)=\E_Q\left[\left(\frac{dP}{dQ}\right)^m\right],
 \qquad D_m(P\mid Q)=\frac1{m-1}\log\mathcal P_m(P\mid Q).
\]

\begin{theorem}[Forward integer-order R\'enyi dominance]\label{thm:renyi}
For every admissible $(E,T)$, integer $m\geq2$, $v>0$, and $\Delta>0$,
\begin{equation}\label{eq:powerdominance}
 \exp\left(\frac{m(m-1)\Delta^2}{2v}\right)\leq M_m(q,z)^T.
\end{equation}
The left side is the full-batch transcript moment. If $T>E$, the inequality
is strict, as is the corresponding forward R\'enyi-divergence comparison.
Thus full batch uniquely minimizes this quantity at the supplied variance.
\end{theorem}

\begin{proof}
Interpret the weights in~\eqref{eq:moment} as the law of
$J\sim\operatorname{Binomial}(m,q)$. Jensen's inequality gives
\[
 M_m(q,z)\geq\exp\left(\frac{\E[J(J-1)]z}{2}\right)
 =\exp\left(\frac{m(m-1)q^2z}{2}\right).
\]
After raising both sides to $T$, the exponent simplifies using $Tq=E$:
\[
 T\frac{m(m-1)q^2}{2}\frac{\Delta^2}{Eqv}
 =\frac{m(m-1)\Delta^2}{2v}.
\]
If $T>E$, then $0<q<1$, every binomial weight is positive, and $J(J-1)$ is
not constant. Strict convexity of the exponential gives strictness. Taking
logarithms and dividing by $m-1>0$ preserves the comparison.
\end{proof}

In particular, any forward order-$m$ R\'enyi budget met by a smaller-batch
point at variance $v$ is also met by full batch at the same variance.
The theorem concerns the forward divergence at integer orders. It does not
establish a forward/reverse comparison, cover noninteger orders, or imply
an ordering of the exact two-sided $(\varepsilon,\delta)$ privacy profiles.

\section{Nonparticipation and zero noise}\label{app:zero-noise}
Recall the zero-noise computation in Appendix~\ref{app:calibration}:
for $\Delta>0$, $\varepsilon\geq0$, $T\in\N_{>0}$, and $q\in[0,1]$,
\begin{equation}\label{eq:zero-noise-profile}
 \AO_S(0,\Delta,q,T,\varepsilon)=1-(1-q)^T.
\end{equation}
The formula includes $q=0$, when the two transcript laws coincide.

For a simple reversal, take $E=1$, $T=2$, and $q=1/2$. The zero-noise profile
is $3/4$, whereas at full batch it is one. At $\varepsilon=0$ and
$\delta=4/5$, the smaller-batch calibrated noise is therefore zero, while
$v_{\mathrm{FB}}>0$. The two total variances are
\[
 V_1(2,S)=S/2,\qquad V_1(1,S)=v_{\mathrm{FB}},
\]
so the smaller batch wins whenever $0\leq S<2v_{\mathrm{FB}}$.

This example is driven by nonparticipation: the differing record is never
selected with probability $1/4$. Its target lies above the boundary
$1-e^{-1}<4/5$, so it does not contradict eventual full-batch dominance in
the subcritical regime. The finite reversal in
Theorem~\ref{thm:strict-interior-counterexample} lies below the corresponding
boundary and therefore does not rely on zero-noise feasibility.

\section{Trade-off separation}\label{app:separation}
Ertan and van Dijk~\cite{ertanvandijk2026} study a related one-epoch Poisson
experiment with $q=1/M$ through the geometry of its $f$-DP trade-off curve.
Their lower bound for noise at a given separation has the same critical scale
$1/\sqrt{2\log M}$. Their main analysis concerns shuffling, with a Poisson
transfer under zero-out adjacency. Our calibration is stated for the
canonical add/remove experiment. At a fixed sampling rate, a zero-contribution
ghost record yields the same reduced scalar Poisson pair; the relevant
distinctions here are the surrounding mechanism assumptions and the
quantities being compared. Their separation bound is not an exact
hockey-stick calibration or a finite-horizon variance ordering.

We use the total-variation convention
\[
 \TV(P,Q)=\sup_A|P(A)-Q(A)|,
\]
where the supremum is over measurable events.
For the canonical transcript pair $P_T,Q_T$, define
\[
 f_T(\alpha)=\inf\{1-\E_{P_T}\psi:
     0\leq\psi\leq1,\ \E_{Q_T}\psi\leq\alpha\},
 \qquad0\leq\alpha\leq1,
\]
where $\psi$ ranges over measurable randomized tests. Let
$\operatorname{sep}(f_T)$ be the maximum Euclidean distance from the graph
of $f_T$ to the random-guessing segment. Since
$0\leq f_T(\alpha)\leq1-\alpha$, orthogonal projection onto that segment gives
\[
 \operatorname{sep}(f_T)
 =\frac1{\sqrt2}\sup_{\alpha\in[0,1]}(1-\alpha-f_T(\alpha))
 =\frac{\TV(P_T,Q_T)}{\sqrt2}.
\]
The last equality follows by optimizing
$\E_{P_T}\psi-\E_{Q_T}\psi$; symmetry of the trade-off curve is not needed.
At $\varepsilon=0$, the two-sided privacy profile is total variation.
Theorem~\ref{thm:accountant-limit} therefore gives the exact separation limit
\[
 \operatorname{sep}(f_T)\longrightarrow
 \frac{1-e^{-\rho\Phi(s)}}{\sqrt2}.
\]
This identifies the critical-window limit of the geometric metric for the
canonical Poisson experiment.

\section{Numerical calibration for the variance illustration}
\label{app:numerical-methods}

Figure~\ref{fig:fixed-participation-variance} uses $\Delta=1$,
$\delta=10^{-5}$, $\varepsilon\in\{1,2,4,8\}$, and
$E\in\{1,10,100,1000\}$. For each $E$, the horizons are the distinct
integers obtained by rounding $19$ logarithmically spaced values from $E$
to $10^6E$ and adding $E\{1,2,10,10^2,\ldots,10^6\}$. There are
$316$ calibrations in total, with the actual sampling rate $q=E/T$.

\paragraph{Privacy quantity and inversion.}
The calculation numerically evaluates the two-sided transcript accountant
$\mathcal A_{E,T}$, not a R\'enyi upper bound or the profile of the average.
For one release under the reference law, write
\[
 L(Z)=\log\bigl(1-q+q e^{Z/\sigma-1/(2\sigma^2)}\bigr),
 \qquad Z\sim\mathcal N(0,1),\qquad
 K(z)=T\log\E e^{zL(Z)}.
\]
The implementation evaluates the two positive-part expectations by
bilateral Laplace inversion: it integrates
$e^{K(z)-(z-1)\varepsilon}/[z(z-1)]$ on a vertical line with
$\Re z>1$ for the forward direction, and
$e^{K(z)+z\varepsilon}/[z(z-1)]$ with $\Re z<0$ for the reverse direction.
Their maximum is calibrated to $\delta$. For $q<1$, when
$\varepsilon\geq-T\log(1-q)$, the reverse profile is exactly zero,
since the product likelihood ratio is at least $(1-q)^T$.
Composition is computed through $K$; it does not allocate $T$ observations
or substitute the sparse-limit formula at large $T$.

\paragraph{Numerical controls and checks.}
Gauss--Legendre quadrature integrates the normal coordinate on
$[-R,R]$, initially with $R=\max\{12,2/\sigma+8\}$.
Calibration compares $1024$ and $2048$ nodes, doubling up to $8192$ if
needed; later passes use $R\geq14$. Small-argument expansions and the
identity $\E[q(e^{Z/\sigma-1/(2\sigma^2)}-1)]=0$ control cancellation
in the one-step moment. Adaptive oscillatory quadrature integrates
frequencies on $[0,\infty)$, with limits of $200$ cycles and subdivisions.
Its requested absolute tolerance, scaled to each final directed profile,
is $2\times10^{-12}$ initially and $5\times10^{-13}$ in later passes.
Brent root finding uses absolute and relative tolerances $2\times10^{-8}$
in $\sigma$; integration warnings and nonfinite values are rejected.

The saved finer-resolution pass reevaluates all $316$ rows with
$\max\{4096,2m\}$ nodes, where $m$ is the row's final calibration
resolution, $R=\max\{16,2/\sigma+8\}$, and frequency tolerance
$5\times10^{-13}$. The largest directed-profile change is
$2.59\times10^{-12}$; the largest residual from $10^{-5}$ is
$1.70\times10^{-11}$, with no integration warnings. A separate
discretized privacy-loss calculation using \texttt{dp-accounting} 0.6.0
checks six moderate-horizon points, halving the mesh from
$\varepsilon/5000$ to at least $\varepsilon/20000$ and further as needed.
Each final comparison differs by less than $2\times10^{-9}$.
These resolution and cross-method checks are numerical diagnostics,
not certified bounds on truncation or discretization error.

Full-batch rows use analytical Gaussian calibration directly. A separate
$80$-digit evaluation of the Gaussian formula checks their stored raw
noise against $\sigma_{E,E}=\sqrt{E v_{\mathrm{FB}}}$, with maximum
relative discrepancy below $10^{-12}$. This checks the noise values,
not merely the normalized endpoint $R_E(1)=1$.

\paragraph{Reproduction files.}
The numerical supplement~\cite{numerical-supplement} is separate from
the earlier r05 transcript-proof snapshot. It contains
\nolinkurl{numerical_batch_sweep.py} (grid, calibration, and checks),
\nolinkurl{plot_batch_sweep.py} (rendering),
\nolinkurl{calibration_data.csv}, the saved resolution and independent
privacy-loss reports, and both figure panels:
\nolinkurl{paper_raw_noise_row.pdf} and
\nolinkurl{paper_variance_row.pdf}.
Its README gives exact commands and dependency versions; the manifest
records source and data hashes. The $\varepsilon=1$ rows were reused from
the earlier sweep; the other $237$ rows were calibrated separately.
Plot reproduction uses only the supplied CSV and does not recalibrate it.

\begin{thebibliography}{9}

\bibitem{raisa2024}
O.~R\"ais\"a, J.~J\"alk\"o, and A.~Honkela.
\newblock Subsampling is not Magic: Why Large Batch Sizes Work for
Differentially Private Stochastic Optimisation.
\newblock In \emph{Proceedings of the 41st International Conference on
Machine Learning}, PMLR 235:41959--41981, 2024.
\newblock \url{https://proceedings.mlr.press/v235/raisa24a.html}.

\bibitem{ballewang2018}
B.~Balle and Y.-X.~Wang.
\newblock Improving the Gaussian Mechanism for Differential Privacy:
Analytical Calibration and Optimal Denoising.
\newblock In \emph{Proceedings of the 35th International Conference on
Machine Learning}, PMLR 80:394--403, 2018.
\newblock \url{https://proceedings.mlr.press/v80/balle18a.html}.

\bibitem{mironov2017}
I.~Mironov.
\newblock R\'enyi Differential Privacy.
\newblock In \emph{30th IEEE Computer Security Foundations Symposium},
pages 263--275, 2017.
\newblock \url{https://doi.org/10.1109/CSF.2017.11}.

\bibitem{mironovtalwarzhang2019}
I.~Mironov, K.~Talwar, and L.~Zhang.
\newblock R\'enyi Differential Privacy of the Sampled Gaussian Mechanism.
\newblock arXiv:1908.10530, 2019.
\newblock \url{https://arxiv.org/abs/1908.10530}.

\bibitem{ertanvandijk2026}
M.~B.~Ertan and M.~van Dijk.
\newblock Fundamental Limitations of Favorable Privacy--Utility Guarantees for
DP-SGD.
\newblock arXiv:2601.10237v3, revised August 9, 2026.
\newblock \url{https://arxiv.org/abs/2601.10237v3}.

\bibitem{kalinin2024}
N.~P.~Kalinin.
\newblock Notes on Sampled Gaussian Mechanism.
\newblock arXiv:2409.04636v1, 2024.
\newblock \url{https://arxiv.org/abs/2409.04636v1}.

\bibitem{rational-supplement}
{\raggedright
Source supplement for fixed-epoch large-batch optimality, revision r05,
September 17, 2026.
\newblock Earlier transcript-result snapshot:
\href{poisson-fixed-epoch-large-batch-optimality-review-data-r05.zip}{%
\nolinkurl{poisson-fixed-epoch-large-batch-optimality-review-data-r05.zip}}.
\newblock Distributed alongside the PDF, with archive SHA-256 in the
accompanying checksum file.
\newblock Includes statement-level formalization coverage, rational bounds,
source manifests, and replay instructions; predates the average-only
extensions in this article.\par}

\bibitem{formal-coverage}
{\raggedright
Statement-level formalization coverage for the current fixed-participation
manuscript, September 30, 2026.
\newblock \href{formal_coverage.md}{\nolinkurl{formal_coverage.md}}.
\newblock Distributed alongside this paper; maps finite and sparse results
to the unchanged r05 Lean snapshot and identifies conventional-only additions.\par}

\bibitem{numerical-supplement}
{\raggedright
Numerical supplement for the fixed-participation variance illustration,
revision r03, September 30, 2026.
\newblock
\href{poisson-fixed-participation-numerical-review-data-r03.zip}{%
\nolinkurl{poisson-fixed-participation-numerical-review-data-r03.zip}}.
\newblock Distributed alongside the PDF with checksums, calibration and
plotting sources, input data, numerical diagnostics, and the figure.\par}

\end{thebibliography}

\end{document}
